% !TeX program = xelatex
% Chinese research manuscript using the Elsevier article class.
\documentclass[preprint,12pt]{elsarticle}
\usepackage[a4paper,margin=25mm]{geometry}
\usepackage[UTF8,fontset=none,scheme=plain]{ctex}
\usepackage{fontspec}
\IfFontExistsTF{Times New Roman}{\setmainfont{Times New Roman}}{\setmainfont{texgyretermes-regular.otf}[BoldFont=texgyretermes-bold.otf,ItalicFont=texgyretermes-italic.otf,BoldItalicFont=texgyretermes-bolditalic.otf]}
\IfFontExistsTF{SimSun}{\setCJKmainfont[AutoFakeBold=2,ItalicFont=SimSun,BoldItalicFont=SimSun]{SimSun}\setCJKmonofont{SimSun}}{\setCJKmainfont[BoldFont=FandolSong-Bold.otf,ItalicFont=FandolSong-Regular.otf,BoldItalicFont=FandolSong-Bold.otf]{FandolSong-Regular.otf}\setCJKmonofont{FandolSong-Regular.otf}}
\usepackage{amsmath,amssymb,unicode-math}
\setmathfont{latinmodern-math.otf}
\usepackage{etoolbox}
\usepackage{graphicx,booktabs,tabularx,array,caption,setspace,microtype,xurl,placeins,longtable,amsthm}
\biboptions{numbers,sort&compress}
\usepackage[hidelinks,unicode]{hyperref}
\hypersetup{pdftitle={音乐游戏活动终分预测：延迟反馈下的跨排名共享与失效诊断},pdfauthor={Tano},pdfsubject={Entertainment Computing 格式中文研究稿}}
\urlstyle{same}
\newcolumntype{L}{>{\raggedright\arraybackslash}X}
\newcolumntype{C}[1]{>{\centering\arraybackslash}p{#1}}
\captionsetup{font=small,labelfont=bf,justification=justified,singlelinecheck=false}
\renewcommand{\figurename}{图}
\renewcommand{\tablename}{表}
\renewcommand{\refname}{参考文献}
\abstracttitle{摘要}
\keywordtitle{\upshape 关键词}
\setstretch{1.15}
\setlength{\parindent}{2em}
\setlength{\parskip}{0.15em}
\setlength{\emergencystretch}{3em}
\hyphenpenalty=10000
\exhyphenpenalty=10000
\setlength{\textfloatsep}{12pt plus 2pt minus 2pt}
\setlength{\intextsep}{12pt plus 2pt minus 2pt}
\renewcommand{\topfraction}{0.92}
\renewcommand{\bottomfraction}{0.9}
\renewcommand{\textfraction}{0.05}
\renewcommand{\floatpagefraction}{0.8}
\setcounter{topnumber}{3}
\setcounter{totalnumber}{4}
\journal{Entertainment Computing}
% Show the manuscript page number without implying submission or acceptance.
\makeatletter
\patchcmd{\pprintMaketitle}{\vskip36pt}{\vskip12pt}{}{\errmessage{Title spacing patch failed}}
\def\ps@pprintTitle{\let\@oddhead\@empty\let\@evenhead\@empty
  \def\@oddfoot{\hfil\thepage\hfil}\let\@evenfoot\@oddfoot}
\makeatother

\newtheorem{proposition}{命题}
\renewcommand{\proofname}{证明}
\newcommand{\upfigure}[4]{\begin{figure}[!htbp]\centering\includegraphics[width=\linewidth,height=0.70\textheight,keepaspectratio]{figures/#1.pdf}\caption{#3}\label{#2}\end{figure}}
% Generated from finalized independent evidence; run scripts/bind_upgrade_evidence.py.
\newcommand{\BaseTest}{7.082}
\newcommand{\BaseVal}{6.672}
\newcommand{\EwmaTest}{6.671}
\newcommand{\EwmaVal}{6.502}
\newcommand{\PairTest}{6.540}
\newcommand{\PairVal}{6.556}
\newcommand{\DifferenceTest}{6.511}
\newcommand{\DlmTest}{6.672}
\newcommand{\PairDelta}{-0.130}
\newcommand{\PairAnchorIncrement}{0.030}
\newcommand{\DlmDelta}{0.001}
\newcommand{\PairTwoTest}{6.5397}
\newcommand{\PairLooLow}{-0.163}
\newcommand{\PairLooHigh}{-0.097}
\newcommand{\PairTopThree}{50.5}
\newcommand{\PairImproved}{11}
\newcommand{\TreeSharedFull}{6.912}
\newcommand{\TreeIndependentFull}{6.312}
\newcommand{\NnSharedFull}{10.949}
\newcommand{\NnHeadsFull}{10.198}
\newcommand{\NnIndependentFull}{9.314}
\newcommand{\CoverageEighty}{93.01}
\newcommand{\WidthEighty}{35.07}
\newcommand{\CoverageNinety}{96.90}
\newcommand{\WidthNinety}{47.64}
\newcommand{\JointEighty}{79.86}
\newcommand{\JointNinety}{90.47}
\newcommand{\ProspectiveRegistered}{2026-11-18 15:02:49.992 UTC}
\newcommand{\ProspectiveTarget}{20}
\newcommand{\ProspectiveEventCount}{0}

% Generated only from sealed evidence.
\newcommand{\RTIndependent}{5.918}
\newcommand{\RSIndependent}{4.977}
\newcommand{\RTQuarter}{5.954}
\newcommand{\RTPool}{6.816}
\newcommand{\RTRisk}{6.017}
\newcommand{\RTGated}{5.918}
\newcommand{\RTPoolDelta}{0.899}
\newcommand{\RTPoolLow}{-0.402}
\newcommand{\RTPoolHigh}{2.012}
\newcommand{\RTRiskDelta}{0.100}
\newcommand{\RTRiskLow}{-0.005}
\newcommand{\RTRiskHigh}{0.225}
\newcommand{\DContemp}{0.4349}
\newcommand{\DLagOne}{0.1549}
\newcommand{\DLagTwo}{-0.0174}
\newcommand{\DQMin}{-0.0010445}
\newcommand{\DCorrLow}{0.1431}
\newcommand{\DCorrHigh}{0.6752}
\newcommand{\DRatio}{7.48}
\newcommand{\DTransfer}{6.953329}
\newcommand{\DTransferLog}{0.008115}
\newcommand{\DTransferI}{5.340927}
\newcommand{\EWidthEighty}{35.20}
\newcommand{\EWinklerEighty}{958.57}
\newcommand{\EWidthNinety}{47.78}
\newcommand{\EWinklerNinety}{1239.23}
\newcommand{\ConfirmTarget}{80}
\newcommand{\ConfirmFloor}{60}
\newcommand{\ConfirmGlobalN}{4}
\newcommand{\ConfirmPairN}{72}
\newcommand{\ConfirmFuture}{0}

\newcommand{\roundfigure}[3]{\IfFileExists{figures/#1.pdf}{\upfigure{#1}{#2}{#3}{round2}}{}}
\begin{document}
\begin{frontmatter}
\title{音乐游戏活动终分预测：\\延迟反馈下的跨排名共享与失效诊断}
\author{Tano}
\begin{abstract}
排名门槛在同一音乐游戏活动中同步增长，终分标签却要等到活动结束后才可用于下一期。本文研究一个具体问题：期初可见的历史误差信号，何时支持跨排名共享，何时应保留独立校正。研究冻结BanG Dream国服2024年1月29日起97期密集采样活动及67/14/16归属，采用预测先封存、结果后释放的活动级时间契约。19档5220个历史测试前缀的活动等权MAPE为\BaseTest\%（未校正）和\EwmaTest\%（独立EWMA）；该测试曾用于探索。进一步只用训练组35期外推活动，预定15/10/10单次外层时间留出，内段逐活动回放后按预定指标从12个候选选择一次，结果保留独立法；末10期独立、完全共享和近期风险共享分别为\RTIndependent\%、\RTPool\%和\RTRisk\%。18项结构消融揭示两条线贡献方向不同。108次重新拟合在总参数匹配和近似前向计算匹配下检验共享结构，优势随终点和历史组成变化。训练残差校准显示同期相关与跨期持续性存在差别；24格条件模拟中潜在状态风险与百分比误差的方向可异；一次固定DLM真实同阶段诊断的MAPE与观测对数平方误差均高于独立法。经验区间的点、完整排名阶段组及整活动覆盖分别审计。本文贡献为可审计时间评测、共享失效的匹配对照及真实量纲机制诊断；基础组合与平方风险公式沿用既有研究。验证反复使用、历史测试探索及回溯元数据限制持续保留，80期确认性采集计划尚未激活，未来成绩为空。
\end{abstract}
\begin{keyword}
音乐游戏 \sep 活动终分预测 \sep 延迟反馈 \sep 时间评测 \sep 多任务学习 \sep 在线误差校正
\end{keyword}
\end{frontmatter}

\section*{Abstract}
{\raggedright\noindent\textbf{Final score prediction in music game events: cross-rank sharing and failure diagnosis under delayed feedback}\par}
We ask when information available at the start of an event supports sharing forecast errors across ranking cutoffs. A frozen cohort contains 97 densely sampled BanG Dream events on the Chinese server, split by event into 67 training, 14 validation and 16 test events. A chronological coordinator seals predictions before releasing completed outcomes. Event macro MAPE over 5220 historical test prefixes is \BaseTest\% without correction and \EwmaTest\% with independent EWMA correction; this test cohort had been explored previously. New designs use only 35 training-period forecast events. A prespecified 15/10/10 design uses one outer temporal diagnostic block and middle-block chronological replay followed by one prespecified selection among 12 candidates, retaining independent correction. On the last ten events, independent, fully pooled and recent-risk corrections score \RTIndependent\%, \RTPool\% and \RTRisk\%, respectively. Eighteen structural ablations separate effects on the two cutoffs. A further 108 neural fits compare sharing at matched total parameter counts and approximately matched forward computation, with outcomes varying by temporal endpoint and history composition. Training residuals link contemporaneous correlation to weak and uncertain temporal persistence. A calibrated 24-regime simulation shows different directions for latent-state risk and percentage error. In one post-hoc diagnostic on real cases from the same phase, a fixed DLM has higher MAPE and observed squared log error than independent correction. Interval audits distinguish point, complete rank--phase group and whole-event coverage. The contribution is an auditable temporal protocol and matched failure diagnosis grounded in real residual scales; forecast-combination and quadratic-risk formulas build on prior work. Reused validation, prior test exploration and retrospective metadata remain limitations. A separate 80-event confirmation plan is unactivated, with no future accuracy reported.

\section{引言}
音乐游戏的限时排名活动持续汇总玩家得分，活动结束时某一排名对应的分数形成进入该排名的门槛。以T1000为例，目标是终场第1000名的得分，同一排名在活动期间可由不同玩家占据。中途预测有助于理解剩余增长和安排投入，但早期前缀、结束前集中增长及活动类型差异会使外推困难。本文研究排名门槛这一汇总量，输入未使用玩家身份或个体行为特征。

图~\ref{fig:upgrade-scene}概括任务。对每个活动，模型接收当前可见的多排名前缀，输出最终分数；历史误差校正的标签来自此前已结束的活动。本期多个排名和多个快照共享环境，独立样本量按活动计。同期线间相关、历史误差的持续性以及任务专属变化分别影响共享的收益。较高的终分相关首先描述活动分数水平，预测价值还取决于已观测前缀解释后留下的变化。

本文把研究问题落实为五项可观测检查：过去已发出的独立与共享误差是否支持下一期组合；原锚点差异配方的收益来自哪条线和哪一项；容量及训练日历控制后共享结构是否仍有优势；真实误差与状态空间模拟相容到什么范围；经验区间和未来活动需要怎样的证据。它们共同检验“同活动相关”向“下一活动可迁移收益”的跨越。

本文的贡献分为三个层次。第一，按活动固定案例、信息时钟和结果释放顺序，保留全部方法、缺失与失败记录。第二，用同支持恒等负控、预定单线消融及容量和计算对照，定位共享效果的来源与取舍。第三，建立训练残差、条件参数集合、模拟与真实迁移之间的可复核连接，并保留新候选选择失败及模型迁移失败。凸组合、EWMA和Kalman滤波采用既有数学；两项命题提供初等机制分析。当前主方法保持独立EWMA，研究候选的额外收益接受固定规则检验。

97期队列已有探索记录，14期验证多次用于选择。新增方法只在训练组内设计和诊断；其末10期仍是回溯训练期外推。本文保留历史暴露、验证复用及回溯元数据三项限制。独立确认性证据由另行激活的新活动队列收集。
\upfigure{upgrade_scene_framework}{fig:upgrade-scene}{多排名终分预测与误差状态的接口。浅侧面区分数据、基础预测和校正模块；它们表达处理层次。前缀截至实际观测时刻，期初历史状态在本期内固定，未知终分沿结束后反馈路径进入后续活动。}{mechanism}

\section{相关工作}
\subsection{累积轨迹预测、延迟反馈与时间评测}
由早期累积轨迹推测长期规模，在在线内容热度预测中已有应用。Szabo与Huberman研究YouTube和Digg访问轨迹及其增长时间尺度\citep{szabo2008popularity}。本文对象具有固定结束时间，预测量为随玩家排序变化的排名门槛，还需处理同活动多排名和批次标签。任务联系在于利用早期轨迹，观测机制与目标含义随场景而异。

Tashman讨论固定或滚动预测起点、重估、窗口和多个测试时期\citep{tashman2000evaluation}；Cerqueira等比较时间序列性能估计方案，并区分平稳与非平稳条件\citep{cerqueira2020evaluation}。延迟在线学习研究给出反馈到达顺序对学习过程的影响\citep{joulani2013delayed}。本文采用活动开始时固定历史、整期预测封存后按严格时间条件释放结果的工程契约，未推导在线遗憾界。活动归属、时间条件成员关系及案例去重则解决本场景的分组与重复计数问题。

\subsection{漂移适应与低维误差校正}
概念漂移研究覆盖变化检测、遗忘和模型更新\citep{gama2014drift}，ADWIN通过在线调整窗口适应数据变化\citep{bifet2007adwin}。本文使用固定系数的几何移动平均\citep{roberts1959ewma}及状态空间控制\citep{kalman1960filter}，没有运行漂移检测。静态误差上升、在线校正下降及反馈延迟敏感性分别提供现象证据，漂移、输入失配和有限样本波动的贡献尚待分离。

两项2026年预印本与低成本适应关系较近：Lade等在多时距自适应共形推断中引入指数加权偏差估计\citep{lade2026bcaci}，Wang等以多遗忘尺度专家及在线聚合适应非平稳预测\citep{wang2026melo}。本文固定每活动、排名及阶段的对数偏差状态，重点分析两线锚点与差异的时间尺度和协方差。基于历史预测误差确定组合权重的研究可追溯至Bates与Granger\citep{bates1969combination}；多任务平均已系统研究任务均值的凸混合与风险收缩\citep{feldman2014averaging}。本轮风险公式与ridge使用这些既有思想，新增设计集中在期初信息集、已发出误差、支持守卫和有限样本诊断。下文经验历史区间提供可靠性诊断，其时间依赖覆盖性质由后续专门研究检验。

\subsection{多任务共享与小样本模型}
联合学习的任务组合影响合作与竞争，Standley等研究任务分组及精度和成本的权衡\citep{standley2020tasks}；PCGrad针对梯度干扰提出优化方法\citep{yu2020pcgrad}。任务关系的正则学习还可显式处理负相关与异常任务\citep{zhang2010relationships}。本文比较共享输出、专属头、相关线分支及参数完全独立的网络。专属头的负结果仍同时受参数量、优化和活动样本数影响，本文未测量梯度冲突。

CatBoost和XGBoost为提取曲线特征后的表格回归提供基线\citep{prokhorenkova2018catboost,chen2016xgboost}；LSTM、TCN及分块Transformer代表不同序列表征\citep{hochreiter1997lstm,bai2018tcn,nie2023patchtst}。表格深度模型的比较需要关注样本规模、特征与预算\citep{grinsztajn2022tabular}。TabPFN、TimesFM及Chronos-2提供不同的上下文或预训练预测路线\citep{hollmann2025tabpfn,das2024timesfm,ansari2025chronos2}，本文把既有实验保留于补充材料，主比较集中于固定基础预测上的校正结构。

\section{数据、预测对象与冻结队列}
\subsection{来源、质量条件与采样边界}
数据来自Bestdori公开的BanG Dream活动、排名分数和元数据接口\citep{bestdori2026data}。设活动为\(e\)，排名为\(r\)，国服实际开始和结束时间为\(S_e,E_e\)。名义预测时刻\(t\)的末次观测为\(u_{e,r,t}\le t\)，当前分数为\(x_{e,r,t}>0\)，核实终分为\(Y_{e,r}>0\)。前缀至少含4个点，末次观测滞后至多60分钟。完整历史曲线至少含5个点，终点位于结束前1分钟至结束后5分钟窗口；终点核实在标签释放后进行。

建模统一限定\(S_e\ge1706504400000\ \mathrm{ms}\)，即北京时间2024年1月29日13时。边界按实际时间比较。基础字典可包含早期角色和卡牌，卡牌按国服发布时间过滤。原存档训练诊断显示，边界前186期的每条曲线点数中位数为31、典型间隔165.5分钟；边界后最初12期分别为335.5和30.0分钟。图~\ref{fig:upgrade-density}保留这一局部诊断。日期边界由数据政策指定，诊断覆盖上述198期历史活动；当前建模队列另为97期。
\upfigure{upgrade_sampling_diagnosis}{fig:upgrade-density}{采样密度的历史存档诊断。观测数、间隔及每天观测频率来自边界前186期与边界后最初12期的训练记录。该描述覆盖局部存档，日期边界沿用预设政策。新旧队列同时改变采样和活动组成，汇总误差各按自己的协议解释。}{data}

\subsection{活动归属、排名覆盖与固定案例}
97期合格活动按实际开始时间排序，冻结为67期训练、14期验证和16期测试。所有排名与协议共用归属；活动ID用于标识和稳定排序。训练开始日期为2024年1月29日至2025年11月30日，验证为2025年12月11日至2026年4月26日，测试为2026年5月6日至10月16日。旧283期队列及结果独立保留。

19档为T10、T20、T30、T40、T50、T100、T200、T300、T400、T500、T1000、T1500、T2000、T3000、T4000、T5000、T10000、T20000及T30000。T1500在训练与验证组无合格曲线，测试覆盖14期242前缀；T10验证和测试各覆盖13期。模型使用相同冻结名单，缺失和质量排除均由数据工件记录。T1500首次预测使用共享排名模型和零本线误差状态；活动结束后，其已释放反馈可供后续活动使用。

\section{时间一致评测与封存流程}
\subsection{固定留出、时间网格与滚动历史}
V1评测25\%、50\%、75\%进度及结束前24、6小时。V2设21个名义条件：开始后1、3、6、12、24、48、72小时；10\%、25\%、50\%、75\%、90\%、95\%进度；结束前48、24、12、6、3、1、0.5、0.25小时。同一活动、排名及物理时刻保存一个案例，多轴条件由成员关系引用。总体指标按唯一案例计算，时间分面保留支持数。

V1和V2的预处理、编码、训练及完整历史统计限于训练组，训练样本历史还须排除本期及开始时未结束的活动。V3保留V2案例与归属，允许此前验证或测试活动结束后用于后续活动。其期初历史为
\begin{equation}\label{eq:history}
\mathcal H_e=\{j:S_j\ge S_{\min},\ E_j<S_e,\ j\text{的结果已按流程释放}\}.
\end{equation}
历史按结束时间、实际开始时间及标识稳定排序。同一期所有时间点和排名共用期初历史池。另一排名前缀裁剪到目标线实际观测时刻，卡牌按国服发布时间过滤。
\begin{table}[!htbp]\centering
\caption{冻结协议的案例数及历史信息条件}\label{tab:protocols}
\small\setlength{\tabcolsep}{4pt}\renewcommand{\arraystretch}{1.15}
\begin{tabularx}{\linewidth}{L r r r L}\toprule
协议范围 & 训练 & 验证 & 测试 & 终分反馈 \\
\midrule
V1两线 & 2390 & 130 & 154 & 不回灌 \\
V2两线 & 2452 & 488 & 560 & 不回灌 \\
V2 19线 & 21984 & 4361 & 5220 & 不回灌 \\
V3两线 & 2452 & 488 & 560 & 已结束活动用于后续 \\
V3 19线 & 21984 & 4361 & 5220 & 已结束活动用于后续 \\
\bottomrule\end{tabularx}
\par\vspace{4pt}\begin{minipage}{\linewidth}\footnotesize 三组活动为67/14/16。V3训练列是父数据集前缀量；32期基础热身后35期形成逐期外推误差，二者角色分列。V3 19线为后续冻结扩展。\end{minipage}
\end{table}


\subsection{暖启动、可信协调器与回退}
暖启动首先用32期训练活动建立基础预测器，其后35期训练活动逐期预测、封存并释放，形成外推残差。再处理14期验证与16期测试，共65期查询。初始32期的拟合误差单独保存，滚动成绩从外推查询计算。图~\ref{fig:upgrade-timeline}显示活动归属和反馈顺序。

可信协调器读取隔离来源并执行式~\eqref{eq:history}，模型只接收已释放历史和盲前缀。V3验证用于选择已登记规则，测试开始前封存选择。测试期间固定规则更新状态，每期预测先封存，再按合法时间释放真实结果。中断运行保留原工件；完全相同配置、源码和依赖才恢复，修订运行使用新目录。
\begin{table}[!htbp]\centering
\caption{逐活动预测、封存与反馈的伪代码}\label{alg:seal}
\small\begin{tabularx}{\linewidth}{r L}\toprule
步骤&可信协调器及模型动作\\\midrule
1&读取冻结活动归属、案例和配置；校验数据、源码、依赖及既有封存指纹。\\
2&按实际开始时间进入活动\(e\)。只释放已经封存预测且结束时间严格早于\(S_e\)的历史结果。\\
3&以期初\(\mathcal H_e\)拟合预处理和基础模型，或读取同指纹缓存；由历史外推预测构造活动阶段残差。\\
4&历史为空时偏差取零；本线支持不足时收缩；配对历史为空时回退本线独立状态；未见排名使用共享基础预测。\\
5&冻结本期状态。逐一接收盲前缀，执行基础解码、校正及共同边界，核对案例覆盖并封存预测与审计记录。\\
6&本期结束并达到释放条件后，由协调器记录结果与释放时刻。下一活动重新执行步骤2。\\
7&独立评分器在预测封存后读答案，输出活动级损失、明细与评分源码指纹。\\\bottomrule
\end{tabularx}\end{table}
\upfigure{upgrade_information_timeline}{fig:upgrade-timeline}{活动级信息顺序。32期基础热身与35期外推误差热身承担不同角色；14期验证和16期测试沿同一固定规则接收合法历史。重叠活动须等待严格结束条件，当前活动全部时间点共用期初池。}{protocol}

\section{基础预测与结构化误差校正}
\subsection{增长目标、基础路线与共同解码}
定义实际进度\(p=(u-S_e)/(E_e-S_e)\)。基础模型以进度调整后的对数增长拟合
\begin{equation}\label{eq:target}
z=\log(Y/x)+\log p,\qquad
\widehat Y^{(0)}=x\exp\!\left[\operatorname{clip}\{f_{\theta_e}(X)-\log p,0,4\}\right].
\end{equation}
边界使预测位于当前分数与其\(\exp(4)\)倍之间；上界服务于计算稳定性。榜单清除异常记录造成下降时，下限可能失配，冻结评测仍保留原案例。

T1000/T2000沿用上下文增强CatBoost400：树深4、学习率0.04、叶节点L2为12、种子42；其余17档为共享CatBoost600：树深5、学习率0.035、L2为12、种子42。前者使用局部增长、形态和乐队、角色、卡牌及历史统计，后者使用对数排名、类型、局部统计和33个归一化形态坐标。该混合基础路线在本轮回放中固定，新增规则未改变基础预测。每个训练活动总权重相同，活动\(e\)的\(m_e\)前缀各取相对权重\(1/m_e\)。

\subsection{活动阶段误差与独立参考}
按实际进度划分早期\([0,.25)\)、中期\([.25,.9)\)及末期\([.9,1)\)。每个已结束活动、排名与阶段汇总封存基础预测的对数误差：
\begin{equation}\label{eq:error}
y_{j,r,q}=\frac{1}{|\mathcal C_{j,r,q}|}\sum_{i\in\mathcal C_{j,r,q}}\log\frac{\widehat Y_i^{(0)}}{Y_i}.
\end{equation}
正值表示偏高。先在活动内求均值，再更新活动序列状态，使反馈按活动计。以首个误差初始化\(v_1=y_1\)，后续\(v_k=(1-\alpha)v_{k-1}+\alpha y_k\)。固定主参考为
\begin{equation}\label{eq:independent}
b^{\mathrm{ind}}_{e,r,q}=C\!\left\{\frac{n_{e,r,q}}{n_{e,r,q}+10}\operatorname{EWMA}_{.3}(y_{1:n,r,q})\right\},\quad C(b)=\operatorname{clip}(b,-.3,.3).
\end{equation}
空历史取零。\(n\)是支持活动数，\(n/(n+10)\)是经验暖启动收缩。固定\(\alpha\)按活动序号衰减，日历半衰期随活动间隔变化。

\subsection{锚点与差异状态}
历史双线组合以T2000为锚点，T1000相对误差为差异。设锚点与配对支持数分别为\(n_a,n_d\)，则
\begin{equation}\label{eq:anchor}
A=\frac{n_a}{n_a+10}\operatorname{EWMA}_{.2}(y_{j,2000,q}),\qquad
D=\frac{n_d}{n_d+5}\operatorname{EWMA}_{.6}(y_{j,1000,q}-y_{j,2000,q}),
\end{equation}
\begin{equation}\label{eq:pair}
b^{\mathrm{pair}}_{2000,q}=C(A),\qquad b^{\mathrm{pair}}_{1000,q}=C(A+D).
\end{equation}
锚点用本线全部合法活动，差异仅用同活动同阶段的配对历史。无配对历史时T1000回退式~\eqref{eq:independent}；锚点无自身历史时使用本线状态。其余17档保持独立参考。全部方法执行共同外部边界
\begin{equation}\label{eq:decode}
\widehat Y=B\{\widehat Y^{(0)}\exp(-b)\},\qquad B(a)=\operatorname{clip}\{a,x,x\exp(4)\}.
\end{equation}
图~\ref{fig:upgrade-anchor}区分坐标、时间尺度与非线性边界。\(A\)的统计含义是指定参照线的偏差，其变化还包含该线的任务因素。
\upfigure{upgrade_anchor_difference}{fig:upgrade-anchor}{锚点、差异与最终校正。两项衰减系数为0.2和0.6，支持收缩常数为10和5；正偏差在解码时扣除。缺配对历史保留本线回退。采用相同滤波、相同支持和相同收缩后，坐标分解满足命题~\ref{prop:identity}的恒等式。}{mechanism}

\section{信息共享的机制与风险条件}
\subsection{终分水平、增长与可预测误差}
设\(L_r=\log Y_r,C_r=\log x_r,G_r=L_r-C_r\)。示例模型\(L_r=M+a_r+\zeta_r,C_r=M+b_r+\nu_r\)中，共同活动尺度\(M\)可使终分高度相关；剩余增长消去\(M\)，其关系由剩余项决定。在共同信息集\(X\)下，
\begin{equation}\label{eq:cov}
\operatorname{Cov}(L_r,L_s)=\operatorname{Cov}\{\mathbb E(L_r\mid X),\mathbb E(L_s\mid X)\}+\mathbb E\{\operatorname{Cov}(L_r,L_s\mid X)\}.
\end{equation}
预测器使用前缀后，可利用的剩余结构须重新估计。再设\(y_e=\theta_e+v_e,\theta_e=F\theta_{e-1}+w_e\)，过程创新跨期独立且独立于旧状态，观测噪声\(v_e\)跨期独立，且与各期过程创新及状态独立（允许同一期跨排名噪声相关），记\(\operatorname{Cov}(v_e)=R\)，则\(\operatorname{Cov}(y_e)=\operatorname{Var}(\theta_e)+R\)，而\(\operatorname{Cov}(y_e,y_{e-1})=F\operatorname{Var}(\theta_{e-1})\)。共同瞬时噪声影响同期关联，其跨期贡献为零。相关矩阵因此按终分、增长、同期误差及滞后误差分别报告。

\subsection{同滤波分解的代数负控}
\begin{proposition}[相同线性滤波的坐标恒等式]\label{prop:identity}
两线采用完全相同的活动历史、顺序、阶段及缺测掩码；\(F\)对两序列使用相同确定性权重，\(s\)为相同乘数，分量相加前保持线性。则
\begin{equation}\label{eq:identity}
sF(y_2)+sF(y_1-y_2)=sF(y_1).
\end{equation}
相加后若执行相同总偏差截断和式~\eqref{eq:decode}，最终预测也相同。
\end{proposition}
\begin{proof}
由线性性，左侧为\(sF(y_2+y_1-y_2)=sF(y_1)\)。共同外部函数作用于相同实数，等式保持。首观测初始化的固定系数EWMA权重为\((1-\alpha)^{n-1}\)及\(\alpha(1-\alpha)^{n-i}\)，满足该条件。
\end{proof}
该初等恒等式用于实现负控。额外单线历史、不同收缩、不同速度或分量内非线性各会改变条件；既有组合同时使用不同速度和收缩，相关消融须逐项固定其余设置。

\subsection{无偏估计的共享平方风险}
\begin{proposition}[固定目标下的池化条件]\label{prop:risk}
对固定目标\(\theta_1,\theta_2\)，无偏估计\(X_1,X_2\)的方差为\(v_1,v_2\)，协方差为\(c\)。令\(d=\theta_2-\theta_1\)，以\(X_w=(1-w)X_1+wX_2\)、\(0\le w\le1\)预测任务1，则相对独立估计的平方风险差为
\begin{equation}\label{eq:risk}
\Delta R(w)=-2w(v_1-c)+w^2(d^2+v_1+v_2-2c).
\end{equation}
两任务都采用\((X_1+X_2)/2\)且等权评分时，平均风险差为\(\{d^2-v_1-v_2+2c\}/4\)。因此完全池化降低平均平方风险的充要条件是
\begin{equation}\label{eq:pooling}
d^2<v_1+v_2-2c.
\end{equation}
\end{proposition}
\begin{proof}
混合估计对任务1的偏差为\(wd\)，方差为\((1-w)^2v_1+w^2v_2+2w(1-w)c\)。加偏差平方、减\(v_1\)，得到式~\eqref{eq:risk}。令\(w=1/2\)，分别对两任务求风险差再取平均，即得式~\eqref{eq:pooling}。半正定协方差矩阵保证\(v_1+v_2-2c\ge0\)。
\end{proof}
该式复述两任务凸平均的基本偏差方差权衡\citep{bates1969combination,feldman2014averaging}：任务差异支付平方偏差，共享减少部分估计噪声。用于滚动状态时，无偏性针对本期潜在状态，方差和协方差须纳入滞后与漂移不确定性。已有偏差、随机任务差异及其相关项需要相应扩展。实际MAPE含指数变换、绝对值、边界及不等案例覆盖，式~\eqref{eq:pooling}的结论范围为平方状态风险；MAPE效果由冻结案例独立评分。

\subsection{差异坐标的协方差与传统DLM}
采用给定载荷与创新假设的二维模型
\begin{equation}\label{eq:state}
\theta_{e,1}=c_e+(1-m)d_e,\quad\theta_{e,2}=c_e-md_e,\quad
c_e=c_{e-1}+\eta_e^c,\quad d_e=d_{e-1}+\eta_e^d.
\end{equation}
过程创新\(\eta_e^c,\eta_e^d\)相互独立且跨期独立，方差分别为\(q_c,q_d\)，并与观测噪声独立。\(m=0\)为参照线仅含共同状态，\(m=.5\)为对称任务差异。定义\(a=c-md\)，原始观测噪声为\(R_\theta=\left(\begin{smallmatrix}r_1&c_v\\c_v&r_2\end{smallmatrix}\right)\)，则锚点差异观测\((y_2,y_1-y_2)\)的噪声为
\begin{equation}\label{eq:noise}
R_{a,d}=\begin{pmatrix}r_2&c_v-r_2\\c_v-r_2&r_1+r_2-2c_v\end{pmatrix},\quad
Q_{a,d}=\begin{pmatrix}q_c+m^2q_d&-mq_d\\-mq_d&q_d\end{pmatrix}.
\end{equation}
独立同方差噪声下，观测交叉项为\(-r\)。分别平滑两坐标隐含删除该关系。完整协方差局部水平DLM保留\(Q,R\)，按Kalman递归更新均值和协方差\citep{kalman1960filter}。真实数据控制仅以训练期外推误差估计参数，首次验证前冻结；缺排名时删除对应观测行，仍按固定规则推进状态。合成实验另固定参数，并设置已知生成参数的oracle，两种身份分别报告。

\subsection{共同支持上的训练诊断}
为减少不同面板组成的影响，终分、半程剩余增长、中期原始外推误差和独立校正误差均取同一34期18线交集。35期外推活动中，活动261缺少合格半程增长前缀；T1500无训练支持。图~\ref{fig:upgrade-correlation}及表~\ref{tab:matchedcorr}按基础路线区分两条上下文线、混合线对及其余共享线对，相关系数描述该共同面板。
\begin{table}[!htbp]\centering
\caption{同一34期18线面板的线对相关中位数}\label{tab:matchedcorr}
\small\setlength{\tabcolsep}{4pt}\renewcommand{\arraystretch}{1.15}
\begin{tabularx}{\linewidth}{L r r r}\toprule
对象 & 两上下文线 & 混合路线线对 & 两共享路线线 \\
\midrule
对数终分 & 0.884 & 0.883 & 0.878 \\
半程剩余对数增长 & 0.313 & 0.472 & 0.531 \\
中期原始外推误差 & 0.283 & 0.192 & 0.227 \\
中期独立校正误差 & 0.130 & 0.233 & 0.242 \\
\bottomrule\end{tabularx}
\par\vspace{4pt}\begin{minipage}{\linewidth}\footnotesize 三列分别有1、32、120个线对。每个系数使用相同34期活动；线对共享活动，相关差区间按活动索引重采样。T1500无训练支持，活动261缺半程增长，均退出共同交集。\end{minipage}
\end{table}

\upfigure{upgrade_matched_correlations}{fig:upgrade-correlation}{同一34期18线面板上的Spearman相关。终分、半程剩余增长、中期原始外推误差与独立校正误差均取相同活动，消除原图67、66、35期支持差异。T1000/T2000使用Context400，其余为共享600树，路线差异由分组明示。}{correlation}

\section{实验设计与主要结果}
\subsection{证据层级、终点与完整菜单}
本轮54种固定方法在评分前封存，包含未校正、独立均值/窗口/EWMA、锚点删项与交换、同滤波负控、13项单因素扰动、替代线对、标量与完整协方差滤波、简单增长法，以及7种反馈策略乘3种主规则。校正方法共享65期查询和基础预测，两项简单增长基线在同一案例上另行计算。本轮主方法固定为\texttt{ewma\_reference}，新评分后保持该身份。

历史数据包内329行记录是2026年11月15日首次汇总的14个\((\text{协议},\text{运行})\)组合，包含方法、种子和集成行。后续实验及本轮54方法另存，其计数口径分别列于补充材料。329行用于保留早期选择轨迹。

主指标为活动等权MAPE。活动集合\(\mathcal E\)内唯一案例为\(\mathcal C_e\)：
\begin{equation}\label{eq:metric}
M=\frac{100}{|\mathcal E|}\sum_{e\in\mathcal E}\frac{1}{|\mathcal C_e|}\sum_{i\in\mathcal C_e}\left|\frac{\widehat Y_i-Y_i}{Y_i}\right|.
\end{equation}
活动内案例等权，缺失使各排名权重仍有差别；另报告活动内排名等权敏感性、case-MAPE、MAE、RMSE和P90。MAPE的权重及非对称性有既有讨论\citep{hyndman2006metrics,demyttenaere2016mape}，对数偏差EWMA属于实用状态估计，目标函数与主评分不同。研究过程中，ChatGPT和Codex辅助代码开发、实验执行与记录核查；成绩由封存预测和独立评分器计算，Tano负责审核。

\subsection{固定基础预测上的校正与结构消融}
表~\ref{tab:main}列主要固定控制，完整54种方法见表~\ref{tab:fullmenu}。5220个测试案例来自16期已探索活动。未校正、独立校正和历史双线组合的MAPE分别为\BaseTest\%、\EwmaTest\%和\PairTest\%。后者相对独立参考差值为\PairDelta\ 个百分点；验证值分别为\BaseVal\%、\EwmaVal\%和\PairVal\%，独立参考保持主方法身份。19线双线组合为\PairTest\%，原T1000/T2000的560案例子面板为\PairTwoTest\%，两种范围分别计算。
\begin{table}[!htbp]\centering
\caption{固定19线基础预测上的主要控制（活动等权MAPE）}\label{tab:main}
\small\setlength{\tabcolsep}{4pt}\renewcommand{\arraystretch}{1.15}
\begin{tabularx}{\linewidth}{L r r r}\toprule
方法 & 验证 / \% & 测试 / \% & 案例P90 / \% \\
\midrule
未校正基础预测 & 6.672 & 7.082 & 17.567 \\
独立EWMA（固定主方法） & 6.502 & 6.671 & 16.760 \\
历史锚点差异组合 & 6.556 & 6.540 & 16.559 \\
仅差异状态 & 6.561 & 6.511 & 16.491 \\
仅锚点状态 & 6.593 & 7.305 & 18.529 \\
交换锚点 & 6.566 & 6.926 & 17.113 \\
同滤波锚点负控 & 6.502 & 6.671 & 16.760 \\
同面板独立负控 & 6.502 & 6.671 & 16.760 \\
5期窗口均值 & 6.464 & 6.697 & 16.859 \\
独立局部水平Kalman & 6.591 & 6.762 & 16.875 \\
完整协方差DLM & 6.510 & 6.672 & 16.750 \\
时间比例增长 & 18.985 & 17.699 & 41.982 \\
\bottomrule\end{tabularx}
\par\vspace{4pt}\begin{minipage}{\linewidth}\footnotesize 验证14期4361案例，测试16期5220案例。完整菜单、延迟通道及敏感性点见表~\ref{tab:fullmenu}。各行均为预先固定控制；本表没有按测试误差重选主方法。\end{minipage}
\end{table}


仅差异消融的测试MAPE为\DifferenceTest\%，完整组合相对它增加\PairAnchorIncrement\ 个百分点，活动配对区间为[−0.070,0.143]，3、5、7活动块区间均跨零。该消融把T2000锚点校正置零，T1000只保留差异项，其余线保持独立参考。完整组合的锚点额外收益在该队列中缺少稳定支持。完整协方差DLM为\DlmTest\%，相对独立参考差值\DlmDelta\ 个百分点，所有重采样区间跨零，二者处于相近精度水平。真实回放的同滤波分解与同面板独立控制在浮点容差内等价，与命题~\ref{prop:identity}一致。上述测试排序用于固定消融描述，主方法仍为独立参考。

图~\ref{fig:upgrade-curves}固定选取按实际开始时间排序第1与第8个测试活动，展示T1000/T2000当前分数、封存预测及核实终分。活动选择与误差值无关；连接线只连接冻结时点。图~\ref{fig:upgrade-ablation}保留删项、换锚、同滤波负控与传统滤波结果，以判断结构而非仅比较最终组合。
\upfigure{upgrade_activity_forecasts}{fig:upgrade-curves}{测试面板中预先指定的两个活动。当前分数、不同校正预测和真实终分在相同冻结时点比较；折线用于连接这些时点，图中没有重建连续密集采样轨迹。两条线的分数尺度分别标示。}{examples}
\upfigure{upgrade_structural_ablations}{fig:upgrade-ablation}{固定基础预测上的结构控制。差值以独立EWMA为对照，负值表示该冻结面板MAPE更低；区间按活动重采样。所有固定方法保留，方法排序承担展示功能，主方法仍为独立参考。}{ablation}

\subsection{既有结构实验的证据位置}
先前的训练规模实验用最早16、32、48、67期重新拟合，并反复评价同一14期488案例验证面板；它同时改变了样本量、日历覆盖和活动组成。既有19线残差网络与手工16格模拟分别检验神经结构和指定生成条件。完整方法、三种子、预算及正负区域保留于补充材料。下节采用训练组内部多终点和匹配预算补充这些控制；各实验的输入与反馈条件分别报告。

\section{训练期外推：共享的可迁移性与失效条件}\label{sec:round2}
\subsection{共同信息边界与研究单位}
新增研究均限定训练67期内部。固定基础流程在32期基础热身后形成35期已发出的外推预测；按实际开始时间将这35期分为前15期建历史、次10期选择、末10期时间外推诊断。该设计只有一个外层末10期时间留出块，内段10期逐活动回放，段末按预定指标选择一次；多轮外层交叉验证未执行。选择块和末段分别有372和352个两线物理案例。每期先按严格\(E_j<S_e\)构造期初池，封存全部当前预测，再由协调器向后释放该活动误差；本期各时点使用相同历史。验证与测试答案未用于新候选选择。名义19档在该训练面板实际观测18档，T1500全缺，冻结缺失保持原样。

12候选风险研究和18项结构消融使用相同35期基础预测；公平训练另取三个训练终点后的12期，它们全部属于上述35期。真实活动并集为35，三个阶段和三个种子保留活动重叠关系。末10期来自已有探索队列，证据身份为回溯训练期外推。

\subsection{已发误差驱动的风险候选与预定选择}
对某排名阶段，独立专家\(I\)使用式~\eqref{eq:independent}。共享专家\(S\)先在同活动的完整排名组内等权平均误差，再使用相同衰减、收缩与边界。主组固定为T1000/T2000。当前校正为\(b_\lambda=I+\lambda(S-I)\)，\(0\le\lambda\le1\)。记已发专家的后续误差\(\epsilon=y-I\)、差值\(h=S-I\)，则
\begin{equation}\label{eq:issuedrisk}
R(\lambda)-R(0)=-2\lambda\mathbb E(\epsilon h)+\lambda^2\mathbb E(h^2).
\end{equation}
这是一维预测组合的标准平方风险展开\citep{bates1969combination,feldman2014averaging}。专家值须取该期实际已发出的版本，标签到达后只计算损失系数；历史专家保持原值。设过去活动等权或几何权重下\(A=\overline{h^2}\)、\(C=\overline{\epsilon h}\)、\(n_{\mathrm{eff}}=(\sum w)^2/\sum w^2\)，固定向独立法收缩的候选为
\begin{equation}\label{eq:lambda}
\widehat\lambda=\operatorname{clip}\!\left\{\frac{C}{A+kA/n_{\mathrm{eff}}},0,1\right\}.
\end{equation}
至少5期且\(A>10^{-8}\)时使用该式，其余回退零。平方目标、经验收缩和MAPE分别评价；当前分数边界也保留。完整候选、系数与删除敏感性见补充材料。

次10期按两线活动等权MAPE只选择一次，独立法以\RSIndependent\%与同面板独立精确并列，按事前顺序选中独立。末10期独立为\RTIndependent\%，完全共享为\RTPool\%，近期风险为\RTRisk\%。完全共享相对独立差\RTPoolDelta\ 个百分点，iid活动区间[\RTPoolLow,\RTPoolHigh]；近期风险差\RTRiskDelta，区间[\RTRiskLow,\RTRiskHigh]，循环块敏感性亦跨零。区间固定已发损失，未计入历史探索和候选搜索。

保守门控检查最近12个完整活动的支持、同期相关、任务差异及差异漂移；末段30个活动阶段门均关闭，两线60个系数全为零，输出因此与独立法相同。风险系数的删除跨度较大，显示有限历史下的权重不稳定。固定1/4共享减少了末段平方对数风险，同时MAPE增至\RTQuarter\%；这项损失分歧在全部候选表中保留。训练选择继续支持独立法，主动共享的稳定额外收益仍待新活动。
\roundfigure{upgrade_adaptive_diagnostics}{fig:round2-adaptive}{训练35期内的风险候选诊断。15/10/10划分、全部候选和封存选择共同呈现；末段配对区间以活动为重抽样单位。门控关闭后的零系数对应独立回退，删除跨度属于敏感性描述。}

\subsection{对称单线消融定位原配方的效应}
固定18项消融分别改变单线替换、锚点方向、滤波速度、收缩和支持。末段原配方在两线面板高于独立参考，仅替换T2000慢锚点时方向相反；T1000为锚的反向替换也保留。表~\ref{tab:round2ablations}给出全部结果。仅T2000的负差属于预定描述，末段没有触发新赢家选择。

同滤波分解与同面板独立输出在浮点精度内一致。锚点与差异的实际历史名单始终相同，完整配对支持消融形成零扰动；该结果限制了本面板对缺线支持效应的识别。平均残差和跨期波动也分开检查：原配方在部分阶段减少平均偏差，同时增大跨期标准差和相邻变化。因而均值校正、逐期稳定性和主MAPE各保留自己的评价对象，玩家行为原因留待额外资料。

\subsection{多终点、参数匹配与计算预算}
公平训练在训练序列32、44、56三个终点冻结历史，各预测随后4期，共12期442个唯一案例。每终点使用最近32期完整历史N32full、同首末时期稀疏16期N16span及最近16期N16recent。前两组比较数量与稀疏组成，后两组比较跨度与组成；输入类别支持随组变化。各结构共用联合训练面板、局部特征、增长目标、预处理及活动与活动内排名权重；共享模型接收两线联合监督，独立分支只接收自身排名行。

共享结构为无偏置\(d\to32\)编码器和两个\(32\to1\)头；独立结构为两个含偏置\(d\to16\)编码器及各自输出头。同单元总参数精确同为\(32d+66\)，实际输入\(d=56\)或57，参数1858或1890。共享宽度、偏置位置及监督来源仍变化，容量计数匹配对应这一明确范围。AdamW和无早停全批次训练固定；三个种子为42、123、2026，主结果对最终分数求均值。

表~\ref{tab:round2fairness}显示，共享60轮在9个终点历史单元中3个点值较低，12期合并的三组配对iid区间均跨零。共享30轮近似匹配独立60轮的前向线性MAC，在9个单元中均较高；其每参数更新次数减半。共享60轮前向MAC约为独立60轮的2.035倍，而标量参数更新总数相同。实际108次拟合、5670次优化器更新和计数定义完整保留于补充材料。这些结果控制容量和计算的不同侧面，结构方向随时间位置与历史组成变化。
\begin{table}[!htbp]\centering
\caption{训练期多终点：总参数匹配与近似前向计算匹配}\label{tab:round2fairness}
\small\setlength{\tabcolsep}{4pt}
\begin{tabularx}{\linewidth}{r L r r r r}\toprule
终点 & 历史组 & 总参数 & 独立60轮 & 共享60轮 & 共享30轮 \\
\midrule
32 & N32full & 1890 & 7.651 & 8.415 & 7.780 \\
32 & N16span & 1858 & 7.301 & 8.146 & 8.706 \\
32 & N16recent & 1890 & 7.225 & 8.043 & 8.199 \\
44 & N32full & 1890 & 6.673 & 7.083 & 8.127 \\
44 & N16span & 1858 & 9.689 & 9.849 & 9.740 \\
44 & N16recent & 1890 & 7.505 & 7.243 & 8.426 \\
56 & N32full & 1890 & 11.834 & 11.510 & 15.238 \\
56 & N16span & 1890 & 12.180 & 12.676 & 15.382 \\
56 & N16recent & 1858 & 12.099 & 11.002 & 13.094 \\
\bottomrule\end{tabularx}
\par\vspace{3pt}\begin{minipage}{\linewidth}\footnotesize 单位为三种子最终分数均值集成的MAPE（\%）；每终点4活动，合计12个查询活动、442个唯一案例。同单元总参数均为32d+66。\end{minipage}
\end{table}

\roundfigure{upgrade_fairness_diagnostics}{fig:round2-fairness}{训练内部的公平结构比较。三个终点各4期，三种子最终分数均值为主；共享60轮对应总参数与每参数更新匹配，共享30轮对应近似前向MAC匹配。9个单元共享查询活动，合计12期。}

\subsection{真实量纲校准、条件模拟与迁移失败}
机制校准限于前25期T1000/T2000中阶段误差，含25对水平值、24个相邻差分及23个差分滞后对。同期Pearson为\DContemp，块长3的描述区间为[\DCorrLow,\DCorrHigh]，自身lag1为\DLagOne/\DLagTwo。前25期相邻查询间隔为9至23天，中位10天；过程方差以查询活动序号计时。基础树随合法历史重训，误差变化也承载基础预测器变化。

白观测噪声与稳定随机游走条件下，差分矩满足\(\Gamma_0=Q+2R\)、\(\operatorname{sym}(\Gamma_1)=-R\)。实测原始\(Q\)最小特征值为\DQMin，原矩与PSD投影分别存档；\(R\)两线对角比为\DRatio。预定似然剖面却约束两线等观测方差，900点中按\(2\Delta\mathrm{NLL}\le10\)保留17个描述性支持点，\(\operatorname{tr}Q/\operatorname{tr}R\)范围为1至45。该带依赖网格、局部水平族和阈值；\(m=0\)及较大的差异过程项均受这些约束。真实玩家的共同与私有状态识别仍缺行为资料和充分的模型适配证据。

取三个支持锚点与八种预定扰动形成24格，每格三种子各100条完整路径；重尾、污染、共同或差异突变、反馈延迟和异方差全部保留。方法与格间使用共同随机数，Monte Carlo单位为完整路径。完全共享与半共享在24格的状态MSE和模拟APE均高于独立；近期风险9格降低、15格增加。固定校准DLM在22格降低潜在状态MSE，在14格降低模拟APE，两个损失排序保留差异；完整方向计数见补充材料。

模拟完成后追加一次固定真实迁移：前25期最优\(Q,R,\mu\)及先验保持不变，仅按合法历史更新中阶段状态。末10期120个共同案例中，独立MAPE为\DTransferI\%，固定DLM为\DTransfer\%，平方对数损失也增加。该诊断在看过模拟后提出，逐期结果先收集于内存、运行末整体封存；其证据限于回溯信息集重放，实际前瞻发布时效另待日志。参数失配、真实状态不可观测和损失目标差异共同留在解释范围，失败结果未触发新网格选择。
\roundfigure{upgrade_calibrated_transfer}{fig:round2-transfer}{真实残差、受限参数集合、条件模拟与事后迁移相连。同期共变与跨期信号分别呈现，支持带依赖预定离散模型族；(c)比较潜在状态MSE与模拟无界APE，(d)报告同阶段真实有界MAPE。事后迁移保留看过模拟后的诊断身份。}


\section{稳健性、尾部误差与区间可靠性}
\subsection{活动差值、重采样与改善集中度}
独立评分先计算相同案例的活动级损失差，再以完整活动进行5000次配对重采样，种子42，并采用3、5、7活动的循环移动块\citep{kunsch1989bootstrap}。图~\ref{fig:upgrade-eventwise}保留全部16期，表~\ref{tab:robustness}列活动区间与块区间。历史组合相对独立参考有\PairImproved\ 期下降，逐一删期差值范围为[\PairLooLow,\PairLooHigh]个百分点，绝对差最大的三期占全部绝对活动差的\PairTopThree\%。重采样固定已发布预测和损失，未重跑在线模型及反馈历史；较长块在16期中只有少量有效片段。区间未作多方法探索和历史选择校正，解释保持探索性\citep{cawley2010selection}。
\begin{table}[!htbp]\centering
\caption{固定已发布损失的活动级配对不确定性}\label{tab:robustness}
\small\setlength{\tabcolsep}{4pt}\renewcommand{\arraystretch}{1.15}
\begin{tabularx}{\linewidth}{L r C{3.1cm} C{3.1cm}}\toprule
比较 & 差值 / 百分点 & 活动配对95\%区间 & 5活动块95\%区间 \\
\midrule
历史组合−独立参考 & -0.130 & [-0.255, -0.009] & [-0.249, -0.024] \\
历史组合−仅差异 & 0.030 & [-0.070, 0.143] & [-0.030, 0.094] \\
DLM−独立参考 & 0.001 & [-0.029, 0.030] & [-0.036, 0.038] \\
\bottomrule\end{tabularx}
\par\vspace{4pt}\begin{minipage}{\linewidth}\footnotesize 5000次重采样，种子42。3和7活动块区间及逐活动损失保存在独立统计summary.json。所有区间为探索性，条件于固定方法和已发布损失，未重跑反馈历史、未校正搜索。\end{minipage}
\end{table}

\upfigure{upgrade_eventwise_improvement}{fig:upgrade-eventwise}{16期活动的配对MAPE差值与改善集中程度。活动按实际开始时间排序，每期权重相同。保留正负差和逐一删期范围；所有值条件于固定封存预测，重采样未重新构造后续历史轨迹。}{events}

\subsection{参数、缺测及反馈延迟}
单因素敏感性包含锚点和差异衰减、支持收缩及最低配对支持，图~\ref{fig:upgrade-sensitivity}同时展示验证和已探索测试。参数在评分前固定，展示后没有采用测试最优点作为选择器。7种反馈策略只作用于三种主规则的残差校正通道，覆盖额外24小时、7天、一次查询延迟，确定性10\%/30\%删减，以及缺少T2000或T1000反馈。缓存基础预测仍按原V3合法历史拟合，额外延迟没有触发基础树重训。因此这些结果描述给定固定基础预测时的校正通道敏感性。24小时场景在当前日程下三种方法的预测均保持一致，未形成有效扰动；7天延迟及丢失呈混合方向，缺少整条线的影响更明显。采集系统缺少历史到达日志，模拟情景使用预设时间规则。
\upfigure{upgrade_sensitivity}{fig:upgrade-sensitivity}{单因素参数敏感性。横轴按预先封存的离散值展示，验证和已探索测试分别标示；同时呈现偏差大小、遗忘速度和支持阈值的变化。所有测试点只作描述，主方法与更新规则保持固定。}{sensitivity}

图~\ref{fig:upgrade-rank-stage}分排名与时间阶段报告差值及活动支持。开始后1小时无合格案例，3小时仅8期8前缀；晚期节点与早期节点的覆盖不同。T1500初始无训练及验证支持，其测试均值混合首次零状态与后续本线反馈。类型和乐队分组各仅1至4期，混合了发生时间、时长、排名覆盖及活动条件，结果用于定位误差。
\upfigure{upgrade_rank_stage}{fig:upgrade-rank-stage}{排名与时间条件的误差及支持。方法间共享冻结案例，每个条件保留活动数；时间轴成员关系用于分面，总体评分仍按物理案例去重。T1500冷启动和缺失条件显式保留。}{rank-stage}

\subsection{百分比误差、绝对尾部与经验区间}
未校正、独立校正和双线组合的case-MAPE为6.929\%、6.464\%、6.342\%，案例P90为17.567\%、16.760\%、16.559\%，RMSE约为546.9万、565.2万、565.0万EP。百分比指标下降伴随平方误差上升，高分排名的大幅绝对误差仍需单列。图~\ref{fig:upgrade-interval}并列展示这些口径。

经验80\%和90\%区间以独立EWMA为中心，历史校准分数取同活动、排名与阶段内最大绝对中心对数误差，并使用\(\lceil(n+1)\ell\rceil\)顺序统计量。支持不足先回退合法历史活动阶段池，再保留显式无界或不可用标记。分布变化下的在线共形推断已有专门理论\citep{gibbs2024online}；本文固定的历史分位数配方以回溯可靠性审计为证据范围。

表~\ref{tab:round2reliability}同时报告三个目标。每水平5220个案例对应878个完整排名阶段组，全部来自16期活动。活动等权逐案例覆盖为\CoverageEighty\%/\CoverageNinety\%，完整组联合覆盖为\JointEighty\%/\JointNinety\%；整活动全部组覆盖仅为0/16和1/16。组联合目标与历史最大误差构造对应；整活动目标要求更多事件同时成立。当前案例全部有有限区间，原始案例均值相对宽度为\WidthEighty\%/\WidthNinety\%，活动等权均值为\EWidthEighty\%/\EWidthNinety\%。两种权重各按自身定义报告。
\begin{table}[!htbp]\centering
\caption{三个覆盖目标与区间宽度}\label{tab:round2reliability}
\small\setlength{\tabcolsep}{4pt}
\begin{tabularx}{\linewidth}{L r r r r}\toprule
名义 / \% & 逐案例 / \% & 完整组联合 [95\%区间] & 整活动成功 & 相对宽度 / \% \\
\midrule
80 & 93.01 & 79.86 [73.82,85.17] & 0/16 & 35.20 \\
90 & 96.90 & 90.47 [85.42,94.47] & 1/16 & 47.78 \\
\bottomrule\end{tabularx}
\par\vspace{3pt}\begin{minipage}{\linewidth}\footnotesize 逐案例、完整组联合及宽度均先在活动内汇总，再活动等权。各水平5220案例、878个完整排名×阶段组、16活动；均为有限区间。重抽样固定已发区间。\end{minipage}
\end{table}


条件审计保留76个逐案例及76个完整组分层。T1000和challenge类型的组覆盖低于总体，类型组仅3期；跨排名回退池的联合覆盖为100\%，相对宽度也明显增大。分母与活动级区间见补充表，单期类型和全成功组的固定损失bootstrap会退化。标准区间评分同时惩罚宽度与两侧越界\citep{gneiting2007proper}，80\%/90\%的活动等权Winkler分数为\EWinklerEighty 万/\EWinklerNinety 万EP；不同名义水平采用不同越界成本。这里未据条件结果调整宽度或方法。
\upfigure{upgrade_tail_interval}{fig:upgrade-interval}{点预测尾部与经验区间的三个目标。百分比误差、绝对RMSE和P90分别评价；逐案例、完整排名阶段组与整活动联合覆盖采用各自分母。覆盖、有限可用率和宽度共同报告，活动数为16。}{interval}

\section{前瞻记录与可执行复现}
前瞻接口采用追加记录：登记配置、接收上游预测提交、封存物理案例、在真实结束后登记可信实际释放、按时间守卫导出历史，再独立评分。该模块管理记录与边界；活动发现、联网采集、基础训练和预测器执行由上游系统承担。运行版本绑定实现源码及依赖，跨批次相同活动、排名和时刻的预测被识别为同一物理案例。正式预测与答案工件隔离。

图~\ref{fig:upgrade-prospective}给出闭环。原登记于\ProspectiveRegistered\ 生效，独立参考为主方法，历史组合和完整协方差DLM为影子对照；最初\ProspectiveTarget\ 期目标保留为流程可行性pilot。新增确认性计划与独立补充保持原登记不变，等待新的激活封存和独立起始队列。唯一主比较为历史组合减独立EWMA的19线活动等权MAPE差，负差支持历史组合、正差支持独立校正；两线机制面板为次要估计目标。

规划方差只取训练35期已发活动差。给定改善0.5个百分点、双侧\(\alpha=.05\)、功效.9、\(\rho=.3\)与上95\%的规划\(\sigma\)，19线模型给出\ConfirmGlobalN\ 期，两线给出\ConfirmPairN\ 期。两线规则只改变少数排名，19线平均稀释了差异方差；该条件化模型输出保留于补充表。工程与稳定性采集下限为\ConfirmFloor\ 期，结合两线次面板并按10向上取整，计划采集\ConfirmTarget\ 期。该下限是工程预算，有限样本推断精度仍取决于活动依赖与分布条件。

主区间预定采用按实际开始时间排序的Bartlett/Newey--West HAC\citep{newey1987hac}，带宽\(L=\min\{n-1,\max(1,\lceil n^{1/3}\rceil)\}\)，并报告iid配对t及循环块敏感性。渐近近似要求弱平稳、有限所需矩与弱依赖；AR(1)功效膨胀用于规划，未来HAC的有限样本功效另待实际数据。零差检验和超过0.5个百分点的实用门槛分别评价。新计划尚未激活，已纳入未来活动\ConfirmFuture\ 期，MAPE为空。

未来按新封存后的实际国服开始时间纳入最初80期，人口清单保留无预测与缺标签活动，缺失原因逐项记录，无替换活动。沿用21时间条件、19排名、至少4点前缀、60分钟观测滞后与15分钟提交滞后上限。可信采集者核实终分来源和实际释放时刻；同一期全部时点共用期初池。流程失败使用事前准时发布的固定当前分数回退，晚补预测标记失效；全部固定人口完整后按预定时点评分。

复现包保存配置、代码、依赖、盲预测、释放日志、独立评分源码与文件哈希。归档模式由封存产物重建数值表、图及LaTeX，重新拟合模式使用新目录。图表脚本读取统计产物，不训练或选择方法。二轮绑定器显式接收仓库根与稿件输出路径；原封存B预测保留，新重放入口使用物理案例别名守卫。独立审计在全部35期验证守卫路径与原预测逐例零差。中断的回放v1保留；CSV类别类型解析修订后以v2重跑全部65期，三项控制的逐案例相对误差上限为0、0及\(1.11\times10^{-16}\)。这类一致性检查验证实现与历史预测的对齐，真实未来精度另由新活动评分提供。
\upfigure{upgrade_prospective_loop}{fig:upgrade-prospective}{前瞻登记、上游预测提交、封存、延迟释放和独立评分的记录闭环。模块按真实释放时刻导出后续可用历史；预测与答案各自封存。当前流程等待合格未来活动，图中没有未来精度数值。}{prospective}

\section{讨论}
\subsection{对游戏预测应用的含义}
期初更新低维偏差状态，可与基础预测分开封存和审计。独立EWMA在历史测试及训练期选择中保持实用参照地位。自适应共享的平方风险系数只依赖已发误差，但有限历史带来的删除敏感性较大；保守门控关闭后与独立法同轨。这项回退避免了本面板的额外共享误差，主动共享的稳定收益仍缺证据。

相关性、持久性和预测损失承担不同作用。同期两线共变支持共享假设的研究动机；训练期滞后信号较弱、原始过程矩含负特征值，以及等噪声约束的失配，共同限制了状态模型解释。受约束模拟保留了负迁移和损失排序变化，真实DLM追加迁移又出现损失增加。应用中应同时显示预测时刻、观测滞后、历史支持与区间宽度，按实际目标审视百分比误差和绝对尾部。

容量匹配消除了一个可量化差异，函数类、宽度与监督共享仍有变化。时间覆盖控制也保留活动组成和类别支持差别。结果适合指导后续结构选择与失败诊断，玩家投入、竞争偏好等行为解释需要玩家层资料或额外设计。

\subsection{证据、元数据与统计限制}
测试16期曾参与早期探索，14期验证被反复使用；本轮追加封存约束后续计算，历史暴露身份继续保留。早期曾查看留出盲前缀的特征分布，该操作接触了留出输入，相关特征探索历史一并披露。独立评分器隔离了本轮模型入口与答案，其工程作用和选择偏差的统计作用分别解释。

元数据采用缺少历史时点档案的回溯快照，按国服发布时间过滤卡牌及严格控制终分时间仍留下历史信息条件重建缺口。采样缺失可能与活动条件相关，T10和T1500的支持差异尤其明显。基础模型的两线/17线混合路线改变输入条件，跨路线相关分组和独立学习曲线均因此单列。首轮训练规模实验的参数与模型数量随结构变化；二轮同单元实现总参数精确匹配，并另给近似前向计算控制，仍保留函数类、更新次数、时期组成和输入维度差异。B、C1与C2的真实活动并集为35期，C2的12期全部来自该集合。

统计区间条件于固定损失，未计入全部候选搜索，也未重建在线反馈轨迹。16期活动较少，时间依赖、非等日历间隔和小类型组限制了精度。首轮手工网格及二轮真实量纲校准模拟均使用受控状态族；二轮保留重尾、突变、延迟和异方差扰动。真实潜在状态缺少直接观测，参数范围条件于离散网格与受限噪声假设；描述性支持带的阈值未经有限样本覆盖校准。结论覆盖单一游戏、服务器和密集采样时期，外部适用性由更多场景与前瞻活动检验。

\section{结论}
本文建立音乐游戏活动终分的活动级时间契约，并将“同期相关是否支持下一期共享”分解为可复核的预测、消融、容量与机制检查。已探索19线测试中，独立EWMA的MAPE为\EwmaTest\%；只使用训练35期的内段逐活动回放、段末预定选择仍保留独立。共享候选、单线效应、多终点训练及条件模拟完整保留，结果显示支持规模、时间位置、噪声失配和损失目标会改变收益方向。

同滤波恒等式与平方风险展开提供已有数学下的初等解释。真实校准的固定DLM在事后迁移中增加误差，使模拟与应用之间的证据缺口具体可见。经验区间的完整组覆盖接近预设水平，整活动与部分条件覆盖仍明显较低，宽度代价需同时报告。历史探索、验证复用及回溯元数据继续限制推断；80期确认性计划等待独立激活和真实新活动。

\FloatBarrier
\section*{数据与代码可用性}
原始公开来源为Bestdori\citep{bestdori2026data}。研究仓库为\href{https://github.com/AhsokaTano26/BanG-DreamActivity-fraction-prediction/tree/research/innovation-20261113}{BanG Dream Activity fraction prediction}。首轮工件根目录为\path{analysis/paper_upgrade_records/20261118}，二轮工件位于\path{analysis/paper_upgrade_round2/20261118}，方法与评分源码、运行清单、全菜单预测、统计表及图表输入均有SHA-256记录。补充材料列明完整冻结数据指纹、历史329行的统计口径与新运行标识；公开分支与本地交付工件的发布状态分别核对。

\section*{生成式人工智能使用声明}
OpenAI的ChatGPT和Codex辅助代码开发、实验执行、记录整理、数值核查、LaTeX编排与语言修改。实际成绩由封存预测和独立评分器计算。署名作者负责审核代码、方法、数值及最终稿件内容，确认投稿版本。工具未列为作者。

\clearpage
\appendix
\newcommand{\PaperMainDocument}{}
% This source is also embedded in main.tex; compile independently with XeLaTeX/BibTeX.
\ifdefined\PaperMainDocument
  \newcommand{\EndSupplementDocument}{}
\else
% !TeX program = xelatex
% Chinese research manuscript using the Elsevier article class.
\documentclass[preprint,12pt]{elsarticle}
\usepackage[a4paper,margin=25mm]{geometry}
\usepackage[UTF8,fontset=none,scheme=plain]{ctex}
\usepackage{fontspec}
\IfFontExistsTF{Times New Roman}{\setmainfont{Times New Roman}}{\setmainfont{texgyretermes-regular.otf}[BoldFont=texgyretermes-bold.otf,ItalicFont=texgyretermes-italic.otf,BoldItalicFont=texgyretermes-bolditalic.otf]}
\IfFontExistsTF{SimSun}{\setCJKmainfont[AutoFakeBold=2,ItalicFont=SimSun,BoldItalicFont=SimSun]{SimSun}\setCJKmonofont{SimSun}}{\setCJKmainfont[BoldFont=FandolSong-Bold.otf,ItalicFont=FandolSong-Regular.otf,BoldItalicFont=FandolSong-Bold.otf]{FandolSong-Regular.otf}\setCJKmonofont{FandolSong-Regular.otf}}
\usepackage{amsmath,amssymb,unicode-math}
\setmathfont{latinmodern-math.otf}
\usepackage{etoolbox}
\usepackage{graphicx,booktabs,tabularx,array,caption,setspace,microtype,xurl,placeins,longtable,amsthm}
\biboptions{numbers,sort&compress}
\usepackage[hidelinks,unicode]{hyperref}
\hypersetup{pdftitle={补充材料：音乐游戏活动终分预测},pdfauthor={Tano},pdfsubject={Entertainment Computing 格式中文研究稿}}
\urlstyle{same}
\newcolumntype{L}{>{\raggedright\arraybackslash}X}
\newcolumntype{C}[1]{>{\centering\arraybackslash}p{#1}}
\captionsetup{font=small,labelfont=bf,justification=justified,singlelinecheck=false}
\renewcommand{\figurename}{图}
\renewcommand{\tablename}{表}
\renewcommand{\refname}{参考文献}
\abstracttitle{摘要}
\keywordtitle{\upshape 关键词}
\setstretch{1.15}
\setlength{\parindent}{2em}
\setlength{\parskip}{0.15em}
\setlength{\emergencystretch}{3em}
\hyphenpenalty=10000
\exhyphenpenalty=10000
\setlength{\textfloatsep}{12pt plus 2pt minus 2pt}
\setlength{\intextsep}{12pt plus 2pt minus 2pt}
\renewcommand{\topfraction}{0.92}
\renewcommand{\bottomfraction}{0.9}
\renewcommand{\textfraction}{0.05}
\renewcommand{\floatpagefraction}{0.8}
\setcounter{topnumber}{3}
\setcounter{totalnumber}{4}
\journal{Entertainment Computing}
% Show the manuscript page number without implying submission or acceptance.
\makeatletter
\patchcmd{\pprintMaketitle}{\vskip36pt}{\vskip12pt}{}{\errmessage{Title spacing patch failed}}
\def\ps@pprintTitle{\let\@oddhead\@empty\let\@evenhead\@empty
  \def\@oddfoot{\hfil\thepage\hfil}\let\@evenfoot\@oddfoot}
\makeatother

\newtheorem{proposition}{命题}
\renewcommand{\proofname}{证明}
\newcommand{\upfigure}[4]{\begin{figure}[!htbp]\centering\includegraphics[width=\linewidth,height=0.70\textheight,keepaspectratio]{figures/#1.pdf}\caption{#3}\label{#2}\end{figure}}
% Generated from finalized independent evidence; run scripts/bind_upgrade_evidence.py.
\newcommand{\BaseTest}{7.082}
\newcommand{\BaseVal}{6.672}
\newcommand{\EwmaTest}{6.671}
\newcommand{\EwmaVal}{6.502}
\newcommand{\PairTest}{6.540}
\newcommand{\PairVal}{6.556}
\newcommand{\DifferenceTest}{6.511}
\newcommand{\DlmTest}{6.672}
\newcommand{\PairDelta}{-0.130}
\newcommand{\PairAnchorIncrement}{0.030}
\newcommand{\DlmDelta}{0.001}
\newcommand{\PairTwoTest}{6.5397}
\newcommand{\PairLooLow}{-0.163}
\newcommand{\PairLooHigh}{-0.097}
\newcommand{\PairTopThree}{50.5}
\newcommand{\PairImproved}{11}
\newcommand{\TreeSharedFull}{6.912}
\newcommand{\TreeIndependentFull}{6.312}
\newcommand{\NnSharedFull}{10.949}
\newcommand{\NnHeadsFull}{10.198}
\newcommand{\NnIndependentFull}{9.314}
\newcommand{\CoverageEighty}{93.01}
\newcommand{\WidthEighty}{35.07}
\newcommand{\CoverageNinety}{96.90}
\newcommand{\WidthNinety}{47.64}
\newcommand{\JointEighty}{79.86}
\newcommand{\JointNinety}{90.47}
\newcommand{\ProspectiveRegistered}{2026-11-18 15:02:49.992 UTC}
\newcommand{\ProspectiveTarget}{20}
\newcommand{\ProspectiveEventCount}{0}

% Generated only from sealed evidence.
\newcommand{\RTIndependent}{5.918}
\newcommand{\RSIndependent}{4.977}
\newcommand{\RTQuarter}{5.954}
\newcommand{\RTPool}{6.816}
\newcommand{\RTRisk}{6.017}
\newcommand{\RTGated}{5.918}
\newcommand{\RTPoolDelta}{0.899}
\newcommand{\RTPoolLow}{-0.402}
\newcommand{\RTPoolHigh}{2.012}
\newcommand{\RTRiskDelta}{0.100}
\newcommand{\RTRiskLow}{-0.005}
\newcommand{\RTRiskHigh}{0.225}
\newcommand{\DContemp}{0.4349}
\newcommand{\DLagOne}{0.1549}
\newcommand{\DLagTwo}{-0.0174}
\newcommand{\DQMin}{-0.0010445}
\newcommand{\DCorrLow}{0.1431}
\newcommand{\DCorrHigh}{0.6752}
\newcommand{\DRatio}{7.48}
\newcommand{\DTransfer}{6.953329}
\newcommand{\DTransferLog}{0.008115}
\newcommand{\DTransferI}{5.340927}
\newcommand{\EWidthEighty}{35.20}
\newcommand{\EWinklerEighty}{958.57}
\newcommand{\EWidthNinety}{47.78}
\newcommand{\EWinklerNinety}{1239.23}
\newcommand{\ConfirmTarget}{80}
\newcommand{\ConfirmFloor}{60}
\newcommand{\ConfirmGlobalN}{4}
\newcommand{\ConfirmPairN}{72}
\newcommand{\ConfirmFuture}{0}


\begin{document}
\begin{frontmatter}
\title{补充材料：音乐游戏活动终分预测}
\author{Tano}
\end{frontmatter}
\appendix
\newcommand{\EndSupplementDocument}{\end{document}}
\fi
\section{补充材料：证据来源、完整菜单与审计范围}\label{sec:supplement}
本材料同时提供正文附录与独立版本，完整保留既有实验并补充训练期诊断。数值表使用原生LaTeX文本。只读绑定器核对封存哈希，并在\path{qa/upgrade_number_bindings.json}与\path{qa/round2_number_bindings.json}记录来源、字段、行条件、单位及舍入。图表的输入、代码与导出指纹由独立manifest记录。主文围绕研究问题解释，完整候选、预算及失败记录在此展开。

\subsection{证据层级与运行索引}
\begin{table}[!htbp]\centering
\caption{研究工件与结论范围}\label{tab:evidencelevels}
\small\begin{tabularx}{\linewidth}{p{3.0cm} L L}\toprule
工件 & 内容 & 解释范围\\\midrule
历史329行汇总 & 2026年11月15日首次菜单，14个协议与运行组合 & 同一方法可含种子和集成行，后续新方法另计\\
\path{fixed_replay_v2} & 65期查询、54种固定规则，19线基础预测逐例对齐 & 已探索测试上的完整机制与敏感性比较\\
\path{two_rank_learning_curve_v2} & 16/32/48/67期训练，12次树、48次神经拟合，56评分单元 & 同一14期488案例验证，局部自身特征和新基础配方\\
\path{matched_train_v3} & 34期18线共同支持相关及重采样 & 训练期描述与相关差诊断\\
\path{synthetic} & 16格、3种子、100序列、4种历史、9方法 & 指定生成模型中的机制与Monte Carlo误差\\
二轮B/C1 & 35期，15/10/10；12候选及18消融 & 训练期回溯外推，共用已发基础预测\\
二轮C2 & 3终点、3历史组、3种子；108拟合 & 12期查询全为35期子集，总参数和近似MAC分别控制\\
二轮D & 前25期校准、24格模拟、末10期事后迁移 & 受限机制诊断，整体运行末封存的迁移身份单列\\
二轮E & 独立重算、76+76条件、功效模型 & 固定损失与探索面板，0个未来活动\\
前瞻登记 & 原20期pilot及新80期独立确认计划 & 新计划待激活，未来精度留空\\\bottomrule
\end{tabularx}\end{table}

固定回放运行标识为\path{dcc493f82e83933d6963e2d39a02e6c624057ad6c0932ad24801efbd252b1f09}；其盲预测根清单SHA-256为\path{6c3a034e9e409c12230cb029f17a3268fec1f051fcce1a1eb34f1ee46390128a}。新拟合学习曲线运行标识为\path{ed9ea52183f47ef4ea566d41d5a52f4ff4661df5385957dae4865a29ae22ba67}，盲预测根清单为\path{45b87d8b9b9123eafde8fea0682b49620d6865e24c9a642032935d17ff1ee195}。评分代码、依赖及输出清单有独立分析标识。根目录保存的v1中断运行和较早学习曲线均独立归档，本文只绑定完成且核验的v2运行。

329行的统计键为\((\text{协议},\text{运行},\text{方法/种子/集成})\)，对应\path{data/supplement_all_329_results.csv}。其中14个协议与运行组合的计数沿用首次菜单；它描述该次归档范围。\path{data/supplement_historical_selections.csv}保留原选中方法和轮次，正文既有神经结果来自\path{data/main_models.csv}。本轮固定菜单的选择身份为\texttt{ewma\_reference}，后续评分没有新增测试最小值选择器。

\subsection{冻结数据指纹与案例成员关系}
\begingroup\footnotesize\setlength{\tabcolsep}{3pt}
\begin{longtable}{p{2.2cm} p{12.8cm}}
\caption{冻结数据完整SHA-256标识}\label{tab:datasetids}\\
\toprule
协议 & 数据指纹 \\
\midrule\endfirsthead
\toprule
协议 & 数据指纹 \\
\midrule\endhead
V1两线 & \path{80e354c3a5340df6724a40ef1df490f92c08007311d55849342de509eaf5ee50} \\
V2两线 & \path{d057beca6597dd28b59fde7da08f665ae74af41d96f940da04268c881864a215} \\
V2 19线 & \path{8ffd781fc9f4ebcad9ea691839777458e5e4c7cbeade58978dca15bce2440ad8} \\
V3两线 & \path{85f5924e0b90deb0e1a30f3aa27d2237640cca4d968d96612dda8fee1d9c72ec} \\
V3 19线 & \path{db4e9b9c84bd64a6de5ee30a090cfac83efbe85c2f6dd899304516e7d031cdf1} \\
\bottomrule
\end{longtable}
\endgroup

所有数据归属按实际国服开始时间确定，阈值为\(1706504400000\)毫秒。基础字典与建模活动分别处理，卡牌按国服发布时间过滤。数据manifest、partitions、excluded和unavailable文件保留质量排除及缺失，模型使用统一案例。多时间轴成员关系仅连接到分面明细，总体损失始终对物理案例去重。

T1500无训练及验证曲线，测试14期242前缀；T10的测试为13期248前缀。T1000测试16期281前缀，T2000为16期279前缀，原两线合计560；19线总体为5220。不同排名覆盖及每期前缀数造成活动内排名权重差别，\texttt{event\_rank\_balanced\_mape\_pct}另将活动内每条可见排名等权。该敏感性字段与正文活动内案例等权口径分别保存。

\section{固定方法、初始化与反馈敏感性}
\subsection{54种方法的完整成绩}
\begingroup\footnotesize\setlength{\tabcolsep}{3pt}
\begin{longtable}{p{6.8cm} r r r r}
\caption{完整54种固定规则及全部反馈策略}\label{tab:fullmenu}\\
\toprule
方法标识 & 验证MAPE & 测试MAPE & case-MAPE & P90 \\
\midrule\endfirsthead
\toprule
方法标识 & 验证MAPE & 测试MAPE & case-MAPE & P90 \\
\midrule\endhead
\path{uncorrected} & 6.672 & 7.082 & 6.929 & 17.567 \\
\path{ewma_reference} & 6.502 & 6.671 & 6.464 & 16.760 \\
\path{legacy_pair} & 6.556 & 6.540 & 6.342 & 16.559 \\
\path{time_proportion} & 18.985 & 17.699 & 17.306 & 41.982 \\
\path{history_growth} & 18.051 & 17.915 & 17.615 & 44.236 \\
\path{residual_mean} & 6.186 & 7.122 & 6.942 & 17.870 \\
\path{window5} & 6.464 & 6.697 & 6.533 & 16.859 \\
\path{window10} & 6.251 & 6.706 & 6.521 & 16.676 \\
\path{ewma_fast} & 6.888 & 7.137 & 6.915 & 18.076 \\
\path{pair_anchor_only} & 6.593 & 7.305 & 7.117 & 18.529 \\
\path{pair_difference_only} & 6.561 & 6.511 & 6.323 & 16.491 \\
\path{pair_swapped_anchor} & 6.566 & 6.926 & 6.711 & 17.113 \\
\path{same_filter_pair} & 6.502 & 6.671 & 6.464 & 16.760 \\
\path{same_panel_independent} & 6.502 & 6.671 & 6.464 & 16.760 \\
\path{kalman_level} & 6.591 & 6.762 & 6.560 & 16.875 \\
\path{kalman_trend} & 7.261 & 7.338 & 7.146 & 18.352 \\
\path{dlm_fullcov} & 6.510 & 6.672 & 6.462 & 16.750 \\
\path{pair_anchor_alpha_0p1} & 6.558 & 6.514 & 6.322 & 16.385 \\
\path{pair_anchor_alpha_0p3} & 6.556 & 6.561 & 6.360 & 16.566 \\
\path{pair_difference_alpha_0p3} & 6.503 & 6.623 & 6.420 & 16.738 \\
\path{pair_difference_alpha_0p9} & 6.603 & 6.535 & 6.336 & 16.563 \\
\path{pair_anchor_shrinkage_0} & 6.559 & 6.564 & 6.364 & 16.646 \\
\path{pair_anchor_shrinkage_5} & 6.557 & 6.551 & 6.352 & 16.613 \\
\path{pair_anchor_shrinkage_20} & 6.555 & 6.526 & 6.330 & 16.539 \\
\path{pair_difference_shrinkage_0} & 6.573 & 6.526 & 6.329 & 16.478 \\
\path{pair_difference_shrinkage_2p5} & 6.563 & 6.533 & 6.336 & 16.491 \\
\path{pair_difference_shrinkage_10} & 6.545 & 6.560 & 6.362 & 16.659 \\
\path{pair_minimum_pairs_3} & 6.556 & 6.540 & 6.342 & 16.559 \\
\path{pair_minimum_pairs_5} & 6.556 & 6.540 & 6.342 & 16.559 \\
\path{pair_minimum_pairs_10} & 6.556 & 6.540 & 6.342 & 16.559 \\
\path{pair_T400_T500} & 6.486 & 6.683 & 6.474 & 16.726 \\
\path{pair_T1000_T3000} & 6.555 & 6.581 & 6.380 & 16.615 \\
\path{pair_T2000_T3000} & 6.493 & 6.671 & 6.466 & 16.751 \\
\path{ewma_reference__delay24h} & 6.502 & 6.671 & 6.464 & 16.760 \\
\path{legacy_pair__delay24h} & 6.556 & 6.540 & 6.342 & 16.559 \\
\path{dlm_fullcov__delay24h} & 6.510 & 6.672 & 6.462 & 16.750 \\
\path{ewma_reference__delay7d} & 6.520 & 6.595 & 6.422 & 16.121 \\
\path{legacy_pair__delay7d} & 6.623 & 6.465 & 6.301 & 16.027 \\
\path{dlm_fullcov__delay7d} & 6.528 & 6.606 & 6.431 & 16.083 \\
\path{ewma_reference__delay_one_query} & 6.520 & 6.633 & 6.464 & 16.388 \\
\path{legacy_pair__delay_one_query} & 6.623 & 6.518 & 6.360 & 16.222 \\
\path{dlm_fullcov__delay_one_query} & 6.528 & 6.645 & 6.474 & 16.345 \\
\path{ewma_reference__drop10pct} & 6.547 & 6.676 & 6.483 & 16.760 \\
\path{legacy_pair__drop10pct} & 6.600 & 6.535 & 6.352 & 16.652 \\
\path{dlm_fullcov__drop10pct} & 6.564 & 6.671 & 6.474 & 16.798 \\
\path{ewma_reference__drop30pct} & 6.411 & 6.713 & 6.530 & 16.980 \\
\path{legacy_pair__drop30pct} & 6.418 & 6.555 & 6.382 & 16.749 \\
\path{dlm_fullcov__drop30pct} & 6.417 & 6.727 & 6.538 & 17.032 \\
\path{ewma_reference__missing_anchor2000} & 6.499 & 6.684 & 6.481 & 16.789 \\
\path{legacy_pair__missing_anchor2000} & 6.499 & 6.684 & 6.481 & 16.789 \\
\path{dlm_fullcov__missing_anchor2000} & 6.469 & 6.968 & 6.764 & 17.585 \\
\path{ewma_reference__missing_target1000} & 6.575 & 7.112 & 6.922 & 18.293 \\
\path{legacy_pair__missing_target1000} & 6.572 & 7.104 & 6.916 & 18.297 \\
\path{dlm_fullcov__missing_target1000} & 6.571 & 7.104 & 6.918 & 18.293 \\
\bottomrule
\end{longtable}
\noindent 全部误差单位为\%。延迟/删减只作用于残差校正反馈，基础预测沿用原V3缓存。每方法测试均为16期5220唯一案例；总体评分未连接时间条件成员关系。\par
\endgroup

\texttt{time\_proportion}使用当前分数除实际进度作为预测；\texttt{history\_growth}对同线合法历史，以进度高斯权重（带宽0.15）先计算每活动的对数剩余增长均值，再对活动等权平均；无同线历史时回退时间比例。两者均执行共同外部分数边界。它们提供简单增长控制，基础预测来源有别于残差校正方法。均值、5期或10期窗口与快EWMA分别改变记忆方式；其余基础树预测保持固定。

仅锚点消融对T1000和T2000都扣除\(A\)；仅差异消融对T1000扣除\(D\)，对T2000偏差设零。交换锚点将T1000作为参照并相应重写差异方向。两种同滤波负控共同采用完整配对面板、\(\alpha=.3\)、收缩常数10，最后才执行总偏差截断。同滤波分解与同面板独立输出的对应比较条件完全相同。

单因素扰动固定锚点\(\alpha\in\{.1,.3\}\)、差异\(\alpha\in\{.3,.9\}\)，锚点收缩\(k\in\{0,5,20\}\)、差异收缩\(k\in\{0,2.5,10\}\)和最低配对支持\(n_{\min}\in\{3,5,10\}\)，逐次只替换该项。替代线对固定为T400/T500、T1000/T3000和T2000/T3000。表中的最低支持扰动在当前留出期可能给出相同预测，其作用是记录实际触发情况。

\subsection{完整协方差DLM的训练期参数}
真实DLM在原两线坐标使用二维随机游走、零初始均值。各阶段以训练期已释放的完整配对外推误差估计参数：令\(C\)为误差水平样本协方差，\(V_\Delta\)为相邻误差差分的样本协方差，设
\begin{equation}
\begin{aligned}
R&=\operatorname{diag}\!\left[\max\{.5\operatorname{diag}(V_\Delta),10^{-6}\}\right],\\
Q&=\operatorname{PSD}_{10^{-6}}\{.1[.5C+.5\operatorname{diag}(C)]\},\\
P_0&=\operatorname{PSD}_{10^{-6}}(C+10^{-6}I).
\end{aligned}
\end{equation}
\(\operatorname{PSD}\)将对称矩阵的特征值下限设为\(10^{-6}\)。配对少于8期时，\(Q=.0001I,R=.01I,P_0=.01I\)。各反馈策略在首次验证前分别固定其训练期参数；后续只更新状态，不重估\(Q,R\)。训练期逐期诊断每次仍只使用该期合法训练历史。

滤波使用预测协方差加\(Q\)、创新和Kalman增益，并用Joseph形式更新协方差。缺排名时删去对应观测行，有效信息不足时保持预测步骤。坐标变换控制把均值、\(P_0,Q,R\)及观测载荷同时变换，测试核对两种坐标的结果等价。该传统方法的结构和参数估计选择均在新评分前固定。

合成DLM的固定参数另为\(Q_{a,d}=\operatorname{diag}(.0002,.002)\)、原始\(R_\theta=.01I\)、原始\(P_\theta=.1I\)，避免与真实训练估参混合。删除交叉项控制只删除变换后\(R\)的非对角项，其余设置相同。oracle使用真实生成\(Q,R\)及初始状态，拥有额外信息，排除于真实方法选择。

\subsection{反馈通道与到达时间假设}
基准V3用\(E_j<S_e\)并核对协调器已释放身份。额外24小时、7天和一次查询延迟仅删减校正器获得的残差行；基础树仍是原V3缓存。10\%和30\%删减由活动、排名、阶段与固定种子42的哈希确定，所有方法使用同一掩码；缺线情景分别删除T2000或T1000的残差反馈。

这些场景衡量给定基础预测的校正通道变化。24小时在本日程下没有改变三种主规则预测；7天和一次查询延迟可使部分损失下降，丢失也呈不同方向。缺少一整条线的反馈更直接检验回退与状态可识别性。历史实际采集时刻缺少完整日志，仿真延迟由设定时间规则生成，基础模型整体延迟重训属于后续独立协议。

\section{共享风险推导与补充诊断}
\subsection{部分共享、已有偏差及时间尺度}
按正文的部分共享平方风险推导，设\(D=d^2+v_1+v_2-2c>0\)。任务1的正向共享降低平方风险当且仅当\(v_1>c\)且\(0<w<2(v_1-c)/D\)，同时\(w\le1\)。受限最优权重为\(\operatorname{clip}\{(v_1-c)/D,0,1\}\)。若\(D=0\)，两估计的差在该目标下几乎处处为零，风险差为零。固定半共享\(w=1/2\)对任务1的条件是\(d^2<3v_1-v_2-2c\)，区别于两任务等权完全池化条件。

已有估计偏差分别为\(b_1,b_2\)时，混合偏差相对项为\([(1-w)b_1+w(b_2+d)]^2-b_1^2\)。随机任务差异与估计误差相关时，还须保留交叉项。这些条件使经验残差相关与风险参数的对应关系保持明确；高终分相关首先描述水平，同期残差相关又混合持久状态和瞬时冲击。

独立噪声方差为\(r\)的稳定EWMA方差为\(r\alpha/(2-\alpha)\)，其有效样本量\((2-\alpha)/\alpha\)在\(\alpha=.2,.3,.6\)时约为9、5.67、2.33。支持活动数\(n\)与有效窗口是两个量。随机游走过程方差\(q\)时，常增益一步状态预测MSE为
\begin{equation}
\frac{q+\alpha^2r}{\alpha(2-\alpha)},\qquad\frac qr=\frac{\alpha^2}{1-\alpha}\quad\text{（最优增益条件）}.
\end{equation}
因此0.2和0.6对应的过程/噪声比约为0.05和0.9，明确表达“锚点慢、差异快”的假设。实际配方还有支持收缩与边界，这些公式用于机制解释。

\subsection{共同支持、相关差及滞后筛查}
\begin{table}[!htbp]\centering
\caption{共同34期153线对的相关中位数差}\label{tab:corrcontrasts}
\small\setlength{\tabcolsep}{4pt}\renewcommand{\arraystretch}{1.15}
\begin{tabularx}{\linewidth}{L r C{4.3cm}}\toprule
比较 & 中位相关差 & 活动配对95\%区间 \\
\midrule
剩余增长−终分 & -0.357 & [-0.570, -0.166] \\
原始误差−终分 & -0.657 & [-0.802, -0.474] \\
校正误差−终分 & -0.642 & [-0.794, -0.457] \\
\bottomrule\end{tabularx}
\par\vspace{4pt}\begin{minipage}{\linewidth}\footnotesize 重采样对所有排名使用同一活动索引，避免把153个线对当独立样本。每次先重估相关再求中位数差；5000次、种子42。块3/5/7及基础路线分组均保留在原CSV，解释为探索性。\end{minipage}
\end{table}

共同面板153个线对的Spearman中位数依次为终分0.878、增长0.521、原始误差0.221、独立校正误差0.237。后两者的全局中位关系略有上升；T1000/T2000从原始0.283降至校正0.130，属于特定线对。正文按对象比较，保留这种异质性。

原外推误差历史含35期，滞后一阶仅34对；本轮34期共同面板另为控制半程增长支持而建立。既有循环位移与全局BY筛查\citep{benjamini2001by}未发现\(q<.05\)关系，离散置换水平、小样本与多重校正共同限制检验力。历史滤波还涉及多阶记忆与收缩，单个一阶相关覆盖其部分统计结构。

\subsection{合成网格、Monte Carlo单位及失败区域}
\upfigure{upgrade_synthetic_sharing}{fig:upgrade-synthetic}{首轮固定合成机制实验。16个生成条件、4种历史长度与9种方法完整保留；每格每长度300条完整Monte Carlo路径，跨格复用底层随机数。状态MSE相对独立EWMA的正差表示负迁移，oracle拥有额外信息。}{synthetic}
固定合成配置摘要为\path{45def3eee6388a994e8dbf22d79ec118f0e83d5c340b6475924a7da7bc9b56ba}。16格均从\(c=d=0\)开始，高斯创新及原始两线噪声跨期独立。各格复用标准正态随机数，不同方法共享路径，四种历史长度来自同一序列。独立Monte Carlo单位是每种子内的完整生成路径，每格每长度共300个；格子和历史长度之间的比较具有设计上的相关性。

\begin{table}[!htbp]\centering
\caption{预登记合成网格的完整方向计数}\label{tab:syntheticcount}
\small\setlength{\tabcolsep}{4pt}\renewcommand{\arraystretch}{1.15}
\begin{tabularx}{\linewidth}{L r r r}\toprule
固定方法 & 网格平均状态MSE & 状态风险下降/上升 & 模拟APE下降/上升 \\
\midrule
零校正 & 0.015699 & 16/48 & 16/48 \\
独立EWMA & 0.003454 & 参考 & 参考 \\
完全共享EWMA & 0.008202 & 32/32 & 32/32 \\
同滤波锚点负控 & 0.003454 & 64组数值等价 & 64组数值等价 \\
既有锚点差异 & 0.003652 & 19/45 & 21/43 \\
独立固定Kalman & 0.002791 & 47/17 & 46/18 \\
删除观测交叉项DLM & 0.003066 & 38/26 & 37/27 \\
完整协方差固定DLM & 0.002852 & 43/21 & 48/16 \\
已知生成参数oracle & 0.002285 & 64/0 & 64/0 \\
\bottomrule\end{tabularx}
\par\vspace{4pt}\begin{minipage}{\linewidth}\footnotesize 16格与4历史长度等权，64组合的权重由设计给定。状态MSE使用统一偏差截断；方向相对独立EWMA，负控按10的负12次方容差判数值等价。oracle拥有真实生成参数，其角色为诊断参照。\end{minipage}
\end{table}

低共同漂移、高差异漂移、低噪声且\(m=0\)时，历史16期的锚点差异状态风险减少0.002100，Monte Carlo区间[−0.002448,−0.001753]。相反，共同漂移0.0004、差异漂移0、噪声0.0064、\(m=0\)且历史32期时增加0.001561，区间[0.001235,0.001887]。载荷从\(m=0\)变为0.5时，锚点差异下降组合数从13/32变为6/32。所有负区域保留，参数未据其重调。

表格与原始记录分别保存截断后的\texttt{state\_mse}、截断前\texttt{raw\_state\_mse}和\texttt{synthetic\_unbounded\_ape\_pct}。模拟APE为\(100\operatorname{mean}|\exp(y_{\mathrm{next}}-b)-1|\)，只共享偏差截断，未加入真实游戏分数边界。全输出包含1728行总体、5184行逐种子及1728行配对结果，来源为\path{synthetic/summary.csv}、\path{per_seed.csv}与\path{paired_comparisons.csv}。

\section{新拟合学习曲线、种子与成本}
\subsection{首轮规模曲线与全部56单元}
每种规模按实际开始时间取最早16、32、48、67期，分别有576、1170、1766、2452个训练前缀，预测同一14期488个两线验证案例。规模、日历覆盖及活动组成共同变化。局部特征、增长目标和边界固定；CatBoost为200轮，共享拟合1个模型、独立拟合2个模型。神经训练60轮，三个种子最终分数均值为主。共享输出、专属头和完全独立的参数及优化预算随结构变化，表格保持原身份。
\begin{table}[!htbp]\centering
\caption{相同488验证案例上的训练规模诊断（MAPE / \%）}\label{tab:learning}
\small\setlength{\tabcolsep}{4pt}\renewcommand{\arraystretch}{1.15}
\begin{tabularx}{\linewidth}{r r r r r r}\toprule
训练活动 & 树共享 & 树独立 & 网络共享 & 网络专属头 & 网络独立 \\
\midrule
16 & 11.227 & 9.838 & 13.128 & 13.275 & 12.749 \\
32 & 8.765 & 8.093 & 12.010 & 11.155 & 10.196 \\
48 & 8.429 & 7.979 & 11.811 & 10.991 & 10.203 \\
67 & 6.912 & 6.312 & 10.949 & 10.198 & 9.314 \\
\bottomrule\end{tabularx}
\par\vspace{4pt}\begin{minipage}{\linewidth}\footnotesize 验证活动固定14期。网络三列为三种子最终分数平均后的评分；树固定种子42。当前局部特征和CatBoost200/神经基础拟合与19档回放的混合基础路线分别解释。\end{minipage}
\end{table}

\upfigure{upgrade_training_size}{fig:upgrade-learning}{首轮重新拟合的规模诊断。所有点反复使用同一14期488案例验证面板；训练规模与日历覆盖共同变化，网络参数与模型数随结构变化。二轮公平对照另在训练内部多终点执行。}{learning}
\begingroup\footnotesize\setlength{\tabcolsep}{3pt}
\begin{longtable}{r l l l r r}
\caption{新拟合学习曲线全部56个评分单元}\label{tab:learningseeds}\\
\toprule
训练期数 & 家族 & 结构 & 种子/集成 & MAPE / \% & P90 / \% \\
\midrule\endfirsthead
\toprule
训练期数 & 家族 & 结构 & 种子/集成 & MAPE / \% & P90 / \% \\
\midrule\endhead
16 & catboost & shared & 42 & 11.227 & 26.375 \\
16 & catboost & independent & 42 & 9.838 & 23.141 \\
16 & neural & shared & 42 & 13.099 & 29.416 \\
16 & neural & shared & 123 & 13.545 & 30.137 \\
16 & neural & shared & 2026 & 14.542 & 31.628 \\
16 & neural & shared & mean & 13.128 & 29.590 \\
16 & neural & rank\_heads & 42 & 14.239 & 31.371 \\
16 & neural & rank\_heads & 123 & 14.830 & 31.936 \\
16 & neural & rank\_heads & 2026 & 15.872 & 32.884 \\
16 & neural & rank\_heads & mean & 13.275 & 29.893 \\
16 & neural & independent & 42 & 12.486 & 28.181 \\
16 & neural & independent & 123 & 13.141 & 28.446 \\
16 & neural & independent & 2026 & 14.145 & 32.337 \\
16 & neural & independent & mean & 12.749 & 28.842 \\
32 & catboost & shared & 42 & 8.765 & 21.060 \\
32 & catboost & independent & 42 & 8.093 & 19.528 \\
32 & neural & shared & 42 & 12.208 & 27.145 \\
32 & neural & shared & 123 & 11.794 & 24.650 \\
32 & neural & shared & 2026 & 13.302 & 28.270 \\
32 & neural & shared & mean & 12.010 & 26.569 \\
32 & neural & rank\_heads & 42 & 12.438 & 27.471 \\
32 & neural & rank\_heads & 123 & 10.794 & 24.491 \\
32 & neural & rank\_heads & 2026 & 12.020 & 26.711 \\
32 & neural & rank\_heads & mean & 11.155 & 24.901 \\
32 & neural & independent & 42 & 9.912 & 23.043 \\
32 & neural & independent & 123 & 10.346 & 23.616 \\
32 & neural & independent & 2026 & 11.407 & 25.249 \\
32 & neural & independent & mean & 10.196 & 23.515 \\
48 & catboost & shared & 42 & 8.429 & 21.129 \\
48 & catboost & independent & 42 & 7.979 & 19.681 \\
48 & neural & shared & 42 & 11.731 & 26.955 \\
48 & neural & shared & 123 & 11.554 & 24.686 \\
48 & neural & shared & 2026 & 13.267 & 28.109 \\
48 & neural & shared & mean & 11.811 & 26.174 \\
48 & neural & rank\_heads & 42 & 12.229 & 26.992 \\
48 & neural & rank\_heads & 123 & 10.776 & 25.102 \\
48 & neural & rank\_heads & 2026 & 11.789 & 26.683 \\
48 & neural & rank\_heads & mean & 10.991 & 25.609 \\
48 & neural & independent & 42 & 9.998 & 23.902 \\
48 & neural & independent & 123 & 10.350 & 24.069 \\
48 & neural & independent & 2026 & 11.168 & 24.806 \\
48 & neural & independent & mean & 10.203 & 23.563 \\
67 & catboost & shared & 42 & 6.912 & 16.563 \\
67 & catboost & independent & 42 & 6.312 & 15.613 \\
67 & neural & shared & 42 & 11.004 & 24.449 \\
67 & neural & shared & 123 & 10.359 & 21.352 \\
67 & neural & shared & 2026 & 12.600 & 27.363 \\
67 & neural & shared & mean & 10.949 & 23.753 \\
67 & neural & rank\_heads & 42 & 11.945 & 24.398 \\
67 & neural & rank\_heads & 123 & 9.683 & 21.354 \\
67 & neural & rank\_heads & 2026 & 11.370 & 25.710 \\
67 & neural & rank\_heads & mean & 10.198 & 22.286 \\
67 & neural & independent & 42 & 8.739 & 20.133 \\
67 & neural & independent & 123 & 9.646 & 21.949 \\
67 & neural & independent & 2026 & 10.816 & 24.150 \\
67 & neural & independent & mean & 9.314 & 20.987 \\
\bottomrule
\end{longtable}
\noindent mean指三个种子的最终解码分数均值，区别于三个误差的平均。全部行共用14期488案例。\par
\endgroup

\begingroup\footnotesize\setlength{\tabcolsep}{3pt}
\begin{longtable}{r l l r r r r}
\caption{实际拟合数、参数量及成本}\label{tab:fitbudget}\\
\toprule
训练期数 & 家族 & 结构 & 拟合数 & 每模型参数 & CPU秒 & 墙钟秒 \\
\midrule\endfirsthead
\toprule
训练期数 & 家族 & 结构 & 拟合数 & 每模型参数 & CPU秒 & 墙钟秒 \\
\midrule\endhead
16 & catboost & shared & 1 & — & 0.528 & 0.592 \\
16 & catboost & independent & 2 & — & 0.987 & 0.978 \\
16 & neural & shared & 3 & 1345 & 1.038 & 1.378 \\
16 & neural & rank\_heads & 3 & 1490 & 0.095 & 0.098 \\
16 & neural & independent & 6 & 1345 & 0.092 & 0.092 \\
32 & catboost & shared & 1 & — & 0.554 & 0.548 \\
32 & catboost & independent & 2 & — & 1.178 & 1.280 \\
32 & neural & shared & 3 & 1345 & 0.100 & 0.100 \\
32 & neural & rank\_heads & 3 & 1490 & 0.130 & 0.131 \\
32 & neural & independent & 6 & 1345 & 0.127 & 0.127 \\
48 & catboost & shared & 1 & — & 0.595 & 0.592 \\
48 & catboost & independent & 2 & — & 1.114 & 1.105 \\
48 & neural & shared & 3 & 1345 & 0.110 & 0.110 \\
48 & neural & rank\_heads & 3 & 1490 & 0.146 & 0.146 \\
48 & neural & independent & 6 & 1345 & 0.134 & 0.135 \\
67 & catboost & shared & 1 & — & 0.599 & 0.594 \\
67 & catboost & independent & 2 & — & 1.132 & 1.119 \\
67 & neural & shared & 3 & 1345 & 0.126 & 0.127 \\
67 & neural & rank\_heads & 3 & 1490 & 0.181 & 0.181 \\
67 & neural & independent & 6 & 1345 & 0.160 & 0.160 \\
\bottomrule
\end{longtable}
\noindent CatBoost各200轮，神经各60轮。网络独立每个种子拟合两个模型，参数列为单模型参数，合计预算随模型数变化。首轮初始化开销计入实际时间；这些小规模CPU成本具有环境依赖。\par
\endgroup

完全独立网络为两条排名各自拟合编码器和输出参数；训练预处理按选定子集拟合，参数张量互不共享。独立回归检查修改T2000监督时，T1000输出保持完全一致。各拟合范围内部按活动总权重相同，独立模型使用本线监督。三种子集成对最终解码分数求平均，所有检查点重新加载后核对封存预测。

网络共享输出、专属头与完全独立结构的单模型参数及拟合数见表~\ref{tab:fitbudget}。相同60轮保留训练时长规则，模型数与参数量随结构变化。该规模诊断不主张总计算量相等；实际墙钟和CPU时间用于复核已执行预算。

\subsection{同验证活动的配对差与解释范围}
\begin{table}[!htbp]\centering
\caption{两线新拟合神经集成的活动配对比较}\label{tab:learningcontrasts}
\small\setlength{\tabcolsep}{4pt}\renewcommand{\arraystretch}{1.15}
\begin{tabularx}{\linewidth}{r L r C{3cm} C{3cm}}\toprule
训练活动 & 比较 & 差值 / 百分点 & 活动配对95\%区间 & 5活动块95\%区间 \\
\midrule
16 & 专属头−共享 & 0.147 & [-0.389, 0.628] & [-0.298, 0.747] \\
16 & 独立−共享 & -0.379 & [-1.519, 0.610] & [-1.170, 0.337] \\
16 & 独立−专属头 & -0.526 & [-1.440, 0.238] & [-1.118, -0.005] \\
32 & 专属头−共享 & -0.855 & [-1.626, -0.140] & [-1.390, -0.221] \\
32 & 独立−共享 & -1.814 & [-2.984, -0.724] & [-2.665, -0.876] \\
32 & 独立−专属头 & -0.959 & [-1.618, -0.361] & [-1.385, -0.541] \\
48 & 专属头−共享 & -0.820 & [-1.553, -0.090] & [-1.267, -0.322] \\
48 & 独立−共享 & -1.608 & [-2.712, -0.591] & [-2.350, -0.747] \\
48 & 独立−专属头 & -0.788 & [-1.434, -0.236] & [-1.193, -0.385] \\
67 & 专属头−共享 & -0.751 & [-1.653, 0.162] & [-1.401, -0.042] \\
67 & 独立−共享 & -1.635 & [-2.694, -0.606] & [-2.399, -0.748] \\
67 & 独立−专属头 & -0.884 & [-1.261, -0.491] & [-1.157, -0.604] \\
\bottomrule\end{tabularx}
\par\vspace{4pt}\begin{minipage}{\linewidth}\footnotesize 均以同一14期验证的活动损失重采样，5000次、种子42。该验证已被多轮使用，区间未校正搜索。参数数量和独立模型数不同，预算见表~\ref{tab:fitbudget}。\end{minipage}
\end{table}

67期训练时，独立网络减专属头的MAPE差为−0.884个百分点，活动配对区间[−1.261,−0.491]。较小训练规模时部分比较区间跨零。每个规模按实际开始时间取最早\(N\)期，增大规模同时延长训练覆盖时期并改变活动组成。学习曲线各点复用同一验证面板，训练子集嵌套，点间相关；参数量、模型数和优化差异仍需进一步匹配。此处的验证结果服务于结构诊断，新测试或生产选择另行登记。

\subsection{既有19线神经残差结构}
既有19线网络在同一基础预测与独立校正上学习残差增量。编码器两层宽度16，输出头宽度8和1；专属头按支持收缩，未见排名用共享输出。相关图使用已结束活动的参考误差，每对至少8期、每线至多3邻居，支持不足回退对数排名距离先验。AdamW学习率.002、权重衰减.05、丢弃率.1、60轮及种子42/123/2026固定，历史不足12期回退。

共享输出、专属头、静态相关线与阶段图注意力的历史测试MAPE为6.716\%、7.240\%、7.213\%、7.218\%。专属头相对共享输出差0.524个百分点，既有活动iid区间[0.261,0.844]。相关线增量差−0.027个百分点，区间跨零。其负结果限定于该轮样本、输入与预算，梯度冲突未被实测。
\upfigure{upgrade_neural_comparison}{fig:upgrade-neural}{既有神经残差结构的匹配比较。活动配对与5活动循环块区间沿用该轮独立统计，历史测试已被探索。参数规模与优化仍有变化，比较范围限于已记录结构。}{neural}

\subsection{历史静态基线与预训练模型}
\begin{table}[!htbp]\centering
\caption{既有V2两线静态实验的描述性汇总}\label{tab:historicalstatic}
\small\setlength{\tabcolsep}{4pt}\renewcommand{\arraystretch}{1.15}
\begin{tabularx}{\linewidth}{L r r r}\toprule
方法 & 训练 / \% & 验证 / \% & 测试 / \% \\
\midrule
排名共享CatBoost & 3.611 & 6.430 & 12.755 \\
检索残差模型 & 3.602 & 5.355 & 23.141 \\
双时钟增长网络 & 5.196 & 9.299 & 12.577 \\
进度增长网络 & 5.402 & 9.605 & 12.398 \\
后续模型组合 & 2.277 & 4.658 & 14.062 \\
TabPFNv2上下文增强 & 0.931 & 5.979 & 12.747 \\
TabPFNv2基础配方 & 2.246 & 6.598 & 13.074 \\
多时距进度辅助监督 & 5.587 & 9.262 & 12.129 \\
Chronos-2（零样本） & — & 33.842 & 31.131 \\
TimesFM2.5（零样本） & — & 38.057 & 35.125 \\
双线TCN三种子均值 & 3.000 & 5.560 & 14.392 \\
阶段路由组合 & 2.203 & 4.852 & 13.259 \\
\bottomrule\end{tabularx}
\par\vspace{4pt}\begin{minipage}{\linewidth}\footnotesize 训练列依方法包含训练拟合或上下文重代入；零样本方法留空。输入、权重版本和预算随历史配方变化。该表与V3回放具有不同信息条件。\end{minipage}
\end{table}

历史TabPFN的训练查询包含在训练上下文中，训练列属于上下文重代入；Chronos-2和TimesFM2.5为零样本推理，训练列留空。TimesFM的2024方法论文与实验2.5检查点分别记录，Chronos-2按2025预印本引用。各历史配方的输入表达、权重版本、原选中方法及指纹见原运行清单。跨协议或跨排名总分用于描述各自任务，同案例和同信息条件的增量另以匹配比较判断。

\section{区间目标、尾部与前瞻接口审计}
\subsection{经验区间与联合覆盖}
\begin{table}[!htbp]\centering
\caption{三种覆盖目标的区别}\label{tab:intervaltargets}
\small\setlength{\tabcolsep}{4pt}\renewcommand{\arraystretch}{1.15}
\begin{tabularx}{\linewidth}{r r r r r r}\toprule
名义 / \% & 组数 & 活动数 & 组内联合 / \% & 逐案例 / \% & 整活动联合 / \% \\
\midrule
80 & 878 & 16 & 79.86 & 93.01 & 0.00 \\
90 & 878 & 16 & 90.47 & 96.90 & 6.25 \\
\bottomrule\end{tabularx}
\par\vspace{4pt}\begin{minipage}{\linewidth}\footnotesize 校准目标是同活动、排名和阶段内的全部前缀。整活动联合要求该活动所有排名和阶段均覆盖，是更强的额外诊断；80\%和90\%各有0/16、1/16期整活动全部覆盖。各目标的块3/5/7区间另存完整CSV。\end{minipage}
\end{table}

每组分数为同活动、排名和阶段内最大的绝对中心对数误差。因此名义水平对应组内全部冻结前缀的联合覆盖。逐案例覆盖较高，覆盖对象较弱；整活动全部排名与阶段同时覆盖更强，当前构造未针对该目标校准。可用率、无界回退和有限支持宽度独立报告，避免用无界区间掩盖支持不足。

支持不足先回退到合法历史的活动阶段最大误差池。T1500可在没有自身历史时借用该池，其\texttt{own\_support}与实际采用的\texttt{support}各自保存。当前5220案例均为有限支持，80\%水平采用支持数最小4、90\%最小9。组内联合覆盖接近名义的回溯现象仍处于单一已探索面板；时间依赖及活动变化使交换性假设需要额外检验。块区间、逐排名/阶段/类型/乐队可靠性保留于完整统计表。

\section{训练35期的候选、消融与公平预算}\label{sec:round2supp}
\subsection{全部候选及支持守卫}
固定专家\(I\)为本线\(\alpha=.3,k=10\)独立EWMA；\(S\)在完整排名组内等权平均后采用相同滤波与收缩，内部偏差边界均为\(\pm.3\)。无完整组历史时\(S=I\)。同面板独立控制把\(I\)也限定到完整组活动，区分池化与支持改变。冻结两线面板中两种独立输出精确一致，原外部独立专家与本轮固定独立专家在11536案例的最大分数差也为零，两者来源身份各自保留。

风险系数只使用已发专家的已释放误差。每活动排名阶段内先对物理案例计算\(a_j=\operatorname{mean}h^2\)与\(c_j=\operatorname{mean}(\epsilon h)\)，再按活动权重汇总。全历史风险用\(k=4/12\)；近期风险用衰减.85、\(k=8\)。向零ridge收缩先于截断，避免改变估计量。逐活动删一后重新估计系数的最小至最大跨度仅描述敏感性。

门控最近最多12个完整活动，要求至少8对、原始误差Pearson\(\rho\ge.15\)、\(\mathrm{RMS}(y_1-y_2)/(\mathrm{RMS}(y_1)+\mathrm{RMS}(y_2)+10^{-8})\le.5\)，差分后的相应比例至多.75；常量序列相关按零处理。全部阈值固定。选择块30个活动阶段中仅1个开启，诊断块全部关闭。门控前风险系数及删除范围继续保存，避免将实际零系数与开关前估计混淆。
\begin{table}[!htbp]\centering
\caption{训练内滚动选择：全部12个固定共享候选}\label{tab:round2adaptive}
\small\setlength{\tabcolsep}{4pt}
\begin{tabularx}{\linewidth}{L r r r}\toprule
候选 & 选择10期 / \% & 末10期 / \% & 末段log平方风险 \\
\midrule
独立EWMA（预定选择） & 4.977 & 5.918 & 0.010006 \\
固定1/4共享 & 5.247 & 5.954 & 0.009912 \\
固定1/2共享 & 5.600 & 6.109 & 0.010134 \\
两线完全共享 & 6.534 & 6.816 & 0.011522 \\
风险收缩 k=4 & 5.028 & 5.968 & 0.010038 \\
风险收缩 k=12 & 5.015 & 5.956 & 0.010027 \\
近期风险 k=8 & 5.013 & 6.017 & 0.010095 \\
稳健门控 k=8 & 4.998 & 5.918 & 0.010006 \\
同面板独立 & 4.977 & 5.918 & 0.010006 \\
风险规则仅改T1000 & 5.008 & 5.918 & 0.010007 \\
风险规则仅改T2000 & 4.983 & 6.017 & 0.010094 \\
相邻排名风险规则 & 4.985 & 5.925 & 0.009974 \\
\bottomrule\end{tabularx}
\par\vspace{3pt}\begin{minipage}{\linewidth}\footnotesize 两线主面板；选择块372案例，末段352案例；各10活动等权。次10期按候选顺序处理并列，选中独立EWMA。log平方损失与MAPE分别保留。\end{minipage}
\end{table}

\begin{table}[!htbp]\centering
\caption{末10期的共享系数与删除敏感性}\label{tab:round2lambda}
\small\setlength{\tabcolsep}{4pt}
\begin{tabularx}{\linewidth}{L r r r r}\toprule
规则 & 正系数/状态 & 均值 & 最大值 & 最大删除跨度 \\
\midrule
风险收缩 k=4 & 15/60 & 0.033 & 0.398 & 0.625 \\
风险收缩 k=12 & 15/60 & 0.027 & 0.321 & 0.500 \\
近期风险 k=8 & 27/60 & 0.072 & 0.583 & 0.802 \\
稳健门控 k=8 & 0/60 & 0.000 & 0.000 & 0.802 \\
风险规则仅改T1000 & 6/60 & 0.009 & 0.152 & 0.365 \\
风险规则仅改T2000 & 21/60 & 0.063 & 0.583 & 0.802 \\
相邻排名风险规则 & 19/60 & 0.265 & 1.000 & 1.000 \\
\bottomrule\end{tabularx}
\par\vspace{3pt}\begin{minipage}{\linewidth}\footnotesize 60状态为10活动×2排名×3阶段。门控行的删除跨度属于门控前风险估计；实际系数全为零。删除跨度描述有限支持敏感性。\end{minipage}
\end{table}


相邻组按19档顺序固定成对，末尾单线保持独立；未用相关矩阵重排。T1000所在(1000,1500)组因T1500全缺始终回退，T2000的(2000,3000)组有支持，故该候选在主两线上具有非对称信息条件。B共35次固定专家计算、420次候选混合，无基础重训；候选、输出与选择seal完整保留。

\subsection{全部18项结构消融}
各规则共用相同基础预测、冻结案例与共同分数边界。单线消融把另一线留在独立参考；仅差异对T2000偏差设零，与仅替换T1000保留T2000独立的控制分别定义。统一\(\alpha\)、统一收缩、完整支持及逐项取消收缩按固定配置执行。
\begin{table}[!htbp]\centering
\caption{训练末10期：18项预定结构消融}\label{tab:round2ablations}
\small\setlength{\tabcolsep}{4pt}
\begin{tabularx}{\linewidth}{L r r}\toprule
固定规则 & 两线MAPE / \% & 相对独立差 / pp \\
\midrule
未校正 & 6.744 & 0.827 \\
独立EWMA & 5.918 & 0.000 \\
原锚点差异 & 6.388 & 0.471 \\
仅替换T1000 & 6.439 & 0.522 \\
仅替换T2000慢锚点 & 5.867 & -0.051 \\
T1000作锚，仅替换T2000 & 6.551 & 0.633 \\
仅锚点 & 7.035 & 1.117 \\
仅差异 & 6.355 & 0.437 \\
同滤波坐标负控 & 5.918 & 0.000 \\
同面板独立 & 5.918 & 0.000 \\
统一衰减系数 & 5.944 & 0.027 \\
统一收缩 & 6.321 & 0.403 \\
完整配对支持 & 6.388 & 0.471 \\
锚点不收缩 & 6.424 & 0.506 \\
差异不收缩 & 6.527 & 0.609 \\
全部不收缩 & 6.554 & 0.636 \\
锚点衰减改为.3 & 6.396 & 0.479 \\
差异衰减改为.3 & 5.923 & 0.006 \\
\bottomrule\end{tabularx}
\par\vspace{3pt}\begin{minipage}{\linewidth}\footnotesize 各规则共用10期352案例。方法标识对应封存配置；仅替换某一线时，另一线保留独立校正。该表保留预定消融身份，末段未用于重选。\end{minipage}
\end{table}

均值、跨期标准差、相邻期变化和lag1逐排名阶段保留于\path{C/runs/c1_002_scores/cross_event_stability.csv}及\path{event_rank_phase_residuals.csv}。稳定性量基于末10活动，阶段行保持同活动关联。原始误差、独立校正与原配方的变化共同报告，行为原因与统计机制分开解释。

\subsection{多终点历史选择、全部种子与实际预算}
N32full与N16span具有相同首末训练期，三个终点跨度分别297、305.875、341天；N16recent跨度145、153.875、175天。前缀数随冻结缺失变化。稀疏选择同时改变活动类型与组成；跨历史组编码的类别支持和\(d=56/57\)也变化。各单元自己的训练预处理对共享和独立结构完全相同，独立网络只接收所属排名监督。

共享网络无偏置\(d\to32\)编码器加两个带偏置输出头，总参数\(32d+66\)；两独立网络各为带偏置\(d\to16\to1\)，合计同为\(32d+66\)。参数相等按实际张量逐项核对。优化器为AdamW（学习率.002、权重衰减.05），全批次、无早停；活动等权且活动内排名等权。三种子集成按最终分数算术平均，非偏差或对数预测平均。

共享前向线性MAC为\(\mathrm{epochs}\,n(32d+64)\)，独立为\(\mathrm{epochs}\,n(16d+16)\)。共享每例计算两个头，独立每例只计算自己的头。相同60轮组的标量参数更新数相同，共享MAC约2.035倍；共享30轮的MAC比独立60轮多约1.73\%，每参数更新次数减半。CPU计时只覆盖拟合循环，初始化、输入读取、预处理和序列化成本另计。
\begin{table}[!htbp]\centering
\caption{公平训练的实际更新与计算预算}\label{tab:round2budget}
\small\setlength{\tabcolsep}{4pt}
\begin{tabularx}{\linewidth}{L r r r r}\toprule
结构 & 优化器步数 & 标量参数更新 & 前向线性MAC & 拟合CPU秒 \\
\midrule
independent60 & 3240 & 3044520 & 1175604480 & 0.563 \\
shared30 & 810 & 1522260 & 1195960320 & 0.249 \\
shared60 & 1620 & 3044520 & 2391920640 & 0.517 \\
\bottomrule\end{tabularx}
\par\vspace{3pt}\begin{minipage}{\linewidth}\footnotesize MAC只计训练前向线性层；反向传播、激活、偏置、优化器及预处理另有成本。共享30轮比独立60轮多约1.73\%前向MAC，每参数更新次数减半。\end{minipage}
\end{table}

\begingroup\footnotesize\setlength{\tabcolsep}{3pt}
\begin{longtable}{r p{2.5cm} r r r r}
\caption{训练期公平对照的全部27个种子单元}\label{tab:round2fairseeds}\\
\toprule
终点 & 历史组 & 种子 & 独立60轮 & 共享60轮 & 共享30轮 \\
\midrule\endfirsthead
\toprule
终点 & 历史组 & 种子 & 独立60轮 & 共享60轮 & 共享30轮 \\
\midrule\endhead
32 & N32full & 42 & 8.743 & 9.381 & 9.256 \\
32 & N32full & 123 & 7.868 & 8.804 & 8.135 \\
32 & N32full & 2026 & 8.711 & 8.138 & 7.850 \\
32 & N16span & 42 & 8.526 & 8.047 & 9.671 \\
32 & N16span & 123 & 8.001 & 9.307 & 9.803 \\
32 & N16span & 2026 & 7.933 & 9.309 & 9.522 \\
32 & N16recent & 42 & 8.986 & 8.943 & 9.545 \\
32 & N16recent & 123 & 7.502 & 8.814 & 7.595 \\
32 & N16recent & 2026 & 8.363 & 8.009 & 8.771 \\
44 & N32full & 42 & 8.705 & 7.135 & 8.324 \\
44 & N32full & 123 & 6.440 & 7.195 & 8.389 \\
44 & N32full & 2026 & 6.060 & 7.447 & 8.863 \\
44 & N16span & 42 & 10.731 & 10.080 & 9.386 \\
44 & N16span & 123 & 10.098 & 9.878 & 10.010 \\
44 & N16span & 2026 & 10.263 & 10.437 & 11.504 \\
44 & N16recent & 42 & 8.580 & 7.826 & 8.811 \\
44 & N16recent & 123 & 8.443 & 8.019 & 9.399 \\
44 & N16recent & 2026 & 6.694 & 7.073 & 8.590 \\
56 & N32full & 42 & 14.265 & 13.227 & 16.488 \\
56 & N32full & 123 & 11.271 & 11.295 & 14.290 \\
56 & N32full & 2026 & 10.697 & 11.340 & 16.488 \\
56 & N16span & 42 & 13.683 & 14.515 & 16.299 \\
56 & N16span & 123 & 11.920 & 11.828 & 15.111 \\
56 & N16span & 2026 & 11.891 & 13.073 & 16.494 \\
56 & N16recent & 42 & 13.959 & 10.867 & 12.726 \\
56 & N16recent & 123 & 9.662 & 10.589 & 12.674 \\
56 & N16recent & 2026 & 14.750 & 12.371 & 15.090 \\
\bottomrule
\end{longtable}
单位为活动等权MAPE（\%）。27个种子单元涉及108次物理拟合；共享结构各27次，独立结构54次。种子与重复终点维持活动重叠关系。\par
\endgroup


C2按结构、组、终点及种子保存MAE/RMSE/P90与配对活动差。12期汇总的总参数匹配差，独立减共享60轮在三历史组的iid区间均跨零。三个终点每个仅4期，较长块区间标注支持不足；合并12期的长块敏感性亦保留少量长段的精度限制。九单元及三个种子用于条件比较，活动数按12期计。

\section{受限机制参数、校准模拟与真实迁移}\label{sec:round2mechanism}
\subsection{原始矩、参数支持与适配缺口}
校准主面板为\(.25\le p<.9\)的T1000/T2000活动均值，前25期定参数、末10期诊断。差分矩、PSD投影、偏度峰度、JB筛查、方差前后比及滞后关联完整保存。JB为渐近筛查，前段108个检验中10个\(p<.05\)、2个BH\(q<.05\)，后段分别5和1个；样本9至25且有时间依赖，高斯适配的判断保留检验力范围。

中阶段原始\(Q\)含负特征值，\(R\)对角比为\DRatio，显示短序列波动与受限等噪声模型的适配缺口。完整900点网格固定\(m\in\{0,.5,1\}\)等离散值，描述性\(2\Delta\mathrm{NLL}\le10\)保留17点。最优共同过程项落下边界，噪声相关落上边界；支持点\(m=0\)的解释条件包括等观测方差及独立过程创新。真实共同或私有玩家状态的唯一识别仍需更多观测和更一般模型。
\begin{table}[!htbp]\centering
\caption{受限局部水平族的参数支持范围}\label{tab:round2calibration}
\small\setlength{\tabcolsep}{4pt}
\begin{tabularx}{\linewidth}{L r r r r r}\toprule
参数 & 最优 & 支持下限 & 支持上限 & 删一下限 & 删一上限 \\
\midrule
qc & 0.000000 & 0.000000 & 0.000736 & 0.000000 & 0.000000 \\
qd & 0.005888 & 0.002944 & 0.011776 & 0.002944 & 0.005888 \\
r & 0.000736 & 0.000147 & 0.001472 & 0.000736 & 0.000736 \\
rho & 0.500000 & -0.500000 & 0.500000 & 0.500000 & 0.500000 \\
m & 0.000000 & 0.000000 & 0.000000 & 0.000000 & 0.000000 \\
q\_r\_ratio & 4.000000 & 1.000000 & 45.000000 & 2.000000 & 4.000000 \\
\bottomrule\end{tabularx}
\par\vspace{3pt}\begin{minipage}{\linewidth}\footnotesize 900点离散网格；描述性支持条件为2ΔNLL≤10，保留17点。参数限于预定等观测方差与独立过程创新族。删一范围仅使用删除后所选参数；原表附带似然元数据属于完整25期。\end{minipage}
\end{table}


删一参数确由删除后似然选出，原\path{delete_one_profile.csv}附带的nll、support及样本量字段引用完整25期网格。\path{D_delete_one_metadata_audit.json}将其身份澄清为\texttt{original\_full25\_*}；本表删一范围仅使用删除后所选参数值。H另以标量滤波复算完整网格及一个固定删除点，参数选择一致。真实运行前的块接缝修订改为重抽原活动附着的真实差分贡献，未将拼接段的跳变当成真实相邻差分。

\subsection{24格条件模拟的完整失败区域}
低、中、高三个支持锚点的\(\operatorname{tr}Q/\operatorname{tr}R\)为1/4/45，三者均\(m=0\)；中位比锚点按事前顺序选实际支持点，身份与最优点分开。每锚点固定高斯、方差规范t5、5\%方差规范污染、共同突变、差异突变、1期延迟、3期延迟及异方差八种条件。每格三种子各100路径、长度35，末10期评分；共43200条路径方法记录。格间共享底层随机数，独立单位为每格的完整路径。
\begin{table}[!htbp]\centering
\caption{真实量纲校准模拟中的完整方向计数}\label{tab:round2sim}
\small\setlength{\tabcolsep}{4pt}
\begin{tabularx}{\linewidth}{L r r}\toprule
规则 & 状态MSE减少/增加 & 模拟APE减少/增加 \\
\midrule
完全共享 & 0/24 & 0/24 \\
固定半共享 & 0/24 & 0/24 \\
近期风险 & 9/15 & 9/15 \\
校准完整协方差DLM & 22/2 & 14/10 \\
已知基准协方差oracle & 22/2 & 22/2 \\
\bottomrule\end{tabularx}
\par\vspace{3pt}\begin{minipage}{\linewidth}\footnotesize 3个参数锚点×8扰动，共24格；每格3种子各100路径，方法与格间共用底层随机数。独立Monte Carlo单位为完整路径；24格描述设计覆盖。\end{minipage}
\end{table}


完全共享与半共享的24格配对Monte Carlo区间在状态MSE及模拟APE上均为正。近期风险在三个3期延迟格均增加状态风险，在共同突变格均减少相对独立风险。固定校准DLM的潜在状态MSE和模拟APE方向在部分格不同，真实得分边界也未进入模拟APE。oracle知道基准生成\(Q,R\)和初态，仍用固定局部水平递归，没有突变预知或逐时异方差信息；其扰动格负结果一并保留。

输出\texttt{raw\_state\_mse}指外部统一偏差截断之前，I/S专家内部\(\pm.3\)截断仍保留。\texttt{clipped\_fraction}只计额外外部截断；字段原名及原值未覆盖。二轮校准锚点均\(m=0\)，首轮手工网格中的载荷.5条件承担另外的机制范围。

\subsection{模拟后追加的固定DLM真实迁移}
追加计划在看过模拟后提出，再封存固定参数与代码；无额外参数搜索。仅对中阶段末10期120案例应用前25期最优\(Q,R,\mu\)及\(P_0=.1I\)，合法完整配对历史由25增至34，每期全部时点共用期初状态。H独立标量复算固定递归、当前分数边界及三比较规则，MAPE与原输出一致。

独立法为\DTransferI\%，固定DLM为\DTransfer\%，观测平方对数损失为\DTransferLog；其余逐期对照保留于\path{D/posthoc_transfer_run_001/event_comparisons.csv}。该运行逐期预测在内存收集后整体封存，未提供逐期真实公开发布时效证据。其结论覆盖固定信息集下的事后回溯迁移，正式前瞻结果由新的发布记录评价。

\section{条件可靠性、评分与未来样本量}\label{sec:round2reliability}
\subsection{完整组分层与宽度代价}
每名义水平5220案例、878完整组、16活动；同组所有冻结前缀同时覆盖才记组成功。\path{E/reliability_002/interval_conditions.csv}给76行逐案例条件，\path{interval_full_group_conditions.csv}给保持完整组的76行。排名、阶段、实际进度、类型、支持及区间来源全部保留，未据误差筛选活动。
\begin{table}[!htbp]\centering
\caption{预定条件审计的低覆盖与回退实例}\label{tab:round2conditions}
\small\setlength{\tabcolsep}{4pt}
\begin{tabularx}{\linewidth}{L r r r r}\toprule
条件 & 名义 / \% & 活动/组数 & 联合覆盖 / \% & 95\%活动区间 \\
\midrule
T1000 & 80 & 16/47 & 65.62 & [50.00,81.25] \\
T1000 & 90 & 16/47 & 87.50 & [72.92,97.92] \\
challenge类型 & 80 & 3/168 & 59.97 & [51.85,66.67] \\
challenge类型 & 90 & 3/168 & 74.79 & [62.96,84.21] \\
跨排名回退池 & 80 & 4/12 & 100.00 & [100.00,100.00] \\
跨排名回退池 & 90 & 10/27 & 100.00 & [100.00,100.00] \\
\bottomrule\end{tabularx}
\par\vspace{3pt}\begin{minipage}{\linewidth}\footnotesize 完整76个逐案例分层与76个完整组分层保留于原生CSV。条件间活动、排名、阶段及历史支持组成共同变化。回退池高覆盖伴随较宽区间。\end{minipage}
\end{table}


全活动联合0/16与1/16的固定损失bootstrap分别退化为[0,0]及活动重抽样分布；独立同分布活动二项参考另给Clopper--Pearson区间[0,20.5907\%]及[.1581\%,30.2321\%]。该二项参考另加活动独立假设。festival与medley各1期，bootstrap为单点；全成功低支持组的单点同样反映固定样本观察。支持与排名、时间和活动组成相关，条件差异解释保留这些混杂。

跨排名\texttt{all\_event\_phase\_pool}回退的完整组覆盖为100\%，80\%水平有4期12组，90\%有10期27组；其活动等权相对宽度约113\%/163\%，自身排名阶段来源约34\%/44\%。不同水平回退数量不同，较高覆盖与宽度代价共同报告。\texttt{own\_support}=0的11案例来自1期，各层小分母维持原值。

误覆盖率\(\alpha=1-\ell\)的标准Winkler评分\citep{gneiting2007proper}为
\begin{equation}
\operatorname{IS}_{\alpha}(l,u;Y)=(u-l)+\frac2\alpha(l-Y)\mathbf1(Y<l)+\frac2\alpha(Y-u)\mathbf1(Y>u).
\end{equation}
先在活动内平均再活动等权，80\%/90\%分别为\EWinklerEighty 万/\EWinklerNinety 万EP。按实际终分归一化的评分另作描述；标准分数尺度和名义水平的越界代价各自保留。最大组误差校准与逐案例区间评分回答不同目标，本轮没有按评分重新选宽度。

\subsection{规划模型、采集预算与推断条件}
功效方差只取训练35期已发legacy\_pair减独立EWMA活动MAPE差。主19线样本标准差约.079753个百分点，两线约.717920；两线规则只影响少数排名，19线平均降低差异方差。采用双侧零差检验、非中心t和正态iid卡方上95\%\(\sigma\)，再用\((1+\rho)/(1-\rho)\)作AR(1)长期方差膨胀敏感性。
\begin{table}[!htbp]\centering
\caption{训练差异驱动的两线样本量敏感性}\label{tab:round2power}
\small\setlength{\tabcolsep}{4pt}
\begin{tabularx}{\linewidth}{L r r r}\toprule
假设备择改善 / pp & ρ=0 & ρ=.3 & ρ=.6 \\
\midrule
0.25 & 151 & 279 & 597 \\
0.5 & 40 & 72 & 151 \\
1.0 & 12 & 20 & 40 \\
\bottomrule\end{tabularx}
\par\vspace{3pt}\begin{minipage}{\linewidth}\footnotesize 双侧α=.05、功效.9、训练35期正态iid卡方区间的上限σ；非中心t基于独立正态活动差，AR(1)长期方差膨胀仅为规划近似。该表规划零差检验。\end{minipage}
\end{table}


在改善.5个百分点、功效.9及\(\rho=.3\)时，主19线模型输出\ConfirmGlobalN\ 期，两线次面板为\ConfirmPairN\ 期；20期pilot在两线该条件下的模型功效约38.09\%。这些输出基于明确规划假设。60期工程采集下限与两线n72共同驱动按10向上取整的80期计划，主19线指标及双侧\(\alpha=.05\)保持不变。超过实用改善.5个百分点的移位零假设另在\path{power_usefulness_margin_grid.csv}列出，与零差检验分别规划。

未来主HAC区间采用时间排序活动差\(d_t\)：\(\widehat\gamma_h=n^{-1}\sum_{t=h+1}^n(d_t-\bar d)(d_{t-h}-\bar d)\)，\(\widehat V=\widehat\gamma_0+2\sum_{h=1}^L(1-h/(L+1))\widehat\gamma_h\)，并报告\(\bar d\pm z_{.975}\sqrt{\widehat V/n}\)。带宽\(L=\min(n-1,\max(1,\lceil n^{1/3}\rceil))\)固定。Bartlett协方差估计与弱依赖渐近条件来自既有研究\citep{newey1987hac}；零方差轨迹标记无显著性决策。iid配对t及块长1/3/5/7的固定损失敏感性另列，未来有限样本精度仍由真实依赖决定。

原20期pilot登记完整保留。新的\path{confirmation_plan.json}加独立\path{confirmation_plan_addendum.json}共同规定80期未激活队列；单独读取原JSON中的模型n4会遗漏采集补充。主比较、主指标、差值定义及双侧检验、完整活动人口、失败回退、缺失留存和固定报告时点均在激活前封存。当前未来活动为0，未来MAPE为空。


\section{版本绑定、失败记录与独立重建}
输入、源码、配置、依赖与预测分别封存。H独立用标量算术重建B全部22680个候选活动排名阶段状态，\(\lambda\)最大差为\(3.22\times10^{-15}\)；C实际张量参数和5670次优化器更新逐项核对。H的独立临时项目完整重放B35期，并重拟C2一个预定单元，与原预测逐例一致。原C源码快照的导入闭包缺口在可搬运包中补全；全新依赖安装环境仍保留自己的验收范围。

原B backend的物理案例别名守卫由独立适配器补充，原源码和封存预测保留；新入口为\path{integration/replay_guarded.py}，配合\path{sharing_input_guard.py}。H在冻结35期验证新守卫与原输出零差。D的原块接缝修订在真实运行前完成；删一参数附带的full25似然元数据用独立审计文件澄清，原表保留。E的statistics001为中断，reliability001因非中心t极尾NaN另存修订；本文数字统一来自statistics002/reliability002，失败目录随档案保留。

稿件内的一键数值重建显式传入仓库根，避免目录层级推断：
\begingroup\small
\begin{verbatim}
.venv/bin/python \
  <manuscript>/scripts/bind_round2_evidence.py \
  --repo-root <absolute-repository> \
  --paper-root <absolute-manuscript>
\end{verbatim}
\endgroup
该命令先重建原15张表与宏，再绑定本轮11张表与宏，逐项记录数值来源和SHA。它只读已封存分数，不拟合、重新评分或改变选择。图表由封存数值绘制，随后在稿件目录运行\texttt{make all}，使用XeLaTeX、BibTeX及交叉引用刷新生成完整稿与独立补充材料。\texttt{make verify}核对引用、字体及LaTeX警告；逐页渲染和接触表另存QA目录。

项目全套unittest已由PI按单线程串行执行，469项通过。二轮B15、C19、D19、E12及独立守卫5项夹具另计，共70项通过。E相对路径夹具首次从工作目录误调用的失败日志保留，按仓库根契约重跑通过，代码保持原版本。软件测试检验执行边界与实现，预测精度来自独立评分。

\subsection{尚待完成的外部证据}
当前真实未来活动为0，80期新队列等待新的激活封存。完整人口清单、准时预测、原始来源及实际释放日志仍依赖上游采集和独立审计。跨服务器和跨游戏验证、历史元数据时点快照与玩家层行为资料，是解释和推广当前结论所需的新增证据。所有历史分数保留测试探索、验证复用与回溯元数据限制；本轮35期研究亦保留回溯身份。

\ifdefined\PaperMainDocument\else
\FloatBarrier
\begingroup\raggedright\small
\bibliographystyle{elsarticle-num}
\bibliography{references}
\endgroup
\fi
\EndSupplementDocument

\FloatBarrier
\begingroup\raggedright\small
\bibliographystyle{elsarticle-num}
\bibliography{references}
\endgroup
\end{document}
